\chapter{Centers and boundary actions}
\label{mht:ch:centres-actions}

A closed boundary observable gives a multiplication operator.
Bulk reconstruction also requires its commutation with boundary
operations and its action on other boundary conditions. Supported
operations must extend when insertions meet. Each requirement adds
data to the multiplication operator.

The diagonal bimodule gives an algebraic recipient for these data.
Its endomorphisms can be computed explicitly for polynomial algebras
and for the continuous Hamiltonian envelopes of
Chapter~\ref{mht:ch:hamiltonian-coefficients}. A geometric action has
further requirements. On a curve, the choice of admissible actors
can even decide whether a final action exists. The comparison with
the independently defined mixed bulk must therefore specify its
recipient before asserting that it is an equivalence.

\section{The diagonal bimodule}

Let $k$ be a field of characteristic zero and let $A$ be an ordinary
unital associative $k$-algebra. Write $A^e=A\otimes_k A^{\mathrm{op}}$.
The diagonal bimodule is $A$, with its left and right multiplication.
Its derived endomorphism complex is
\[
 Z_{\mathrm{bimod}}(A)=R\operatorname{Hom}_{A^e}(A,A).
\]
This definition specifies both the object and the category of
modules in which its resolutions are formed.

Put $\bar A=A/k1$. The normalized bar resolution has terms
\[
 B_n(A)=A\otimes\bar A^{\otimes n}\otimes A,
 \qquad n\geq0,
\]
in homological degree $n$. Its boundary is the alternating sum of
adjacent multiplications, including the two end actions. The
augmentation $B_0(A)\to A$ is multiplication. 

Each term is a free
$A^e$-module. Inserting $1$ on the left contracts the augmented
complex as a complex of $k$-vector spaces. The normalized formula
uses the class of the former leftmost factor in $\bar A$.
Consequently this bar complex resolves the diagonal.

Applying $\operatorname{Hom}_{A^e}(-,A)$ gives the normalized
Hochschild cochains
\[
 C^n(A,A)=\operatorname{Hom}_k(\bar A^{\otimes n},A).
\]
For a cochain $F$ of arity $n$, its differential is
\begin{align*}
 (bF)(a_1,\ldots,a_{n+1})
 ={}&a_1F(a_2,\ldots,a_{n+1})\\
 &+\sum_{i=1}^n(-1)^i
 F(a_1,\ldots,a_ia_{i+1},\ldots,a_{n+1})\\
 &+(-1)^{n+1}F(a_1,\ldots,a_n)a_{n+1}.
\end{align*}
The formula is evaluated on representatives, with normalized
cochains vanishing when an argument is $1$. Thus
$C^\bullet(A,A)$ computes $Z_{\mathrm{bimod}}(A)$.

A Hochschild $0$-cochain is an element $a\in A$. It is closed
precisely when $a$ is central. Multiplication by $a$ is then a
bimodule endomorphism of the diagonal. For a commutative dg algebra,
a degree-zero element closed under the internal differential gives
the same construction. This accounts for a boundary multiplication
operator without identifying all derived endomorphisms of the
diagonal.

The cochains carry the cup product and insertion operations. An
arity-$q$ insertion has reduced degree $q-1$. Inserting it after
$i$ arguments contributes $(-1)^{i(q-1)}$ in the ungraded case.
The braces sum over ordered nonoverlapping insertions, with these
signs. For dg inputs their internal Koszul signs are included as
well. 

Write $F\circ G=F\{G\}$. Our unshifted bracket convention is
\[
 [F,G]_G=(-1)^{(p-1)(q-1)}F\circ G-G\circ F
 \qquad(F\in C^p,\ G\in C^q).
\]
This is the suspension normalization for which $b=[m,-]_G$ and
$[X,f]_G=X(f)$ for a derivation $X$. The bracket has degree $-1$.
These operations retain more
information than the graded algebra $HH^\bullet(A,A)$.

\section{Polynomial coordinates and the Poisson differential}

Take $A=k[x_1,\ldots,x_m]$. The diagonal ideal in
$A^e=k[x_1,\ldots,x_m,y_1,\ldots,y_m]$ is generated by the regular
sequence $x_i-y_i$. Its Koszul resolution is
\[
 K=A^e\otimes\Lambda(\eta_1,\ldots,\eta_m),
 \qquad d\eta_i=x_i-y_i,
\]
where the $\eta_i$ have homological degree one. After applying
$\operatorname{Hom}_{A^e}(-,A)$ the differential vanishes. The
result is the polyvector algebra
\[
 \mathrm{PV}(A)=A\otimes\Lambda(\partial_1,\ldots,\partial_m),
 \qquad |\partial_i|=1.
\]
Here $\partial_i$ denotes the constant derivation
$\partial/\partial x_i$.

\begin{proposition}\label{mht:prop:polynomial-hkr}
Antisymmetrization defines a quasi-isomorphism
\[
 \operatorname{HKR}:(\mathrm{PV}(A),0)
             \longrightarrow(C^\bullet(A,A),b),
\]
with
\[
 \operatorname{HKR}(fX_1\wedge\cdots\wedge X_n)(a_1,\ldots,a_n)
 =\frac f{n!}\sum_{\sigma\in S_n}\operatorname{sgn}(\sigma)
           \prod_{i=1}^nX_{\sigma(i)}(a_i).
\]
On cohomology it identifies the exterior product with cup product
and the Schouten bracket with the Gerstenhaber bracket.
\end{proposition}

\begin{proof}
Send $\eta_{i_1}\wedge\cdots\wedge\eta_{i_n}$ to the alternating
sum of bar words $[x_{i_{\sigma(1)}}|\cdots|x_{i_{\sigma(n)}}]$.
Internal bar faces cancel in pairs because the coordinates commute.
The end faces give the Koszul boundary. This is a comparison of
free resolutions lifting the identity of $A$.

Restriction of the displayed cochain to that alternating sum is
the corresponding Koszul dual basis vector: the $n!$ terms cancel
the denominator. It therefore induces the computed isomorphism
on cohomology. Products and brackets can be checked on functions
and derivations. The cup classes of derivations anticommute, the
bracket of two derivations is their commutator, and their action
on a function is differentiation. The derivation identities extend
these formulas to all polyvectors.
\end{proof}

The differential in this proposition is zero on polyvectors.
A Poisson tensor introduces a different complex. Fix the Schouten
convention $[\partial_i,x_j]_S=\delta_{ij}$, and put
$d_\pi=[\pi,-]_S$. For the linear Poisson tensor
\[
 \pi=y\,\partial_x\wedge\partial_y
 \quad\text{on }k[x,y],
\]
one has $d_\pi x=-y\partial_y$. The Hochschild differential of the
$0$-cochain $x$ is still zero in the commutative algebra. Hence HKR
from $(\mathrm{PV},d_\pi)$ to the undeformed Hochschild complex is
not a chain map.

There is nevertheless an exact Chevalley description of this
Poisson differential. Let $\mathfrak l$ be finite dimensional, with
$[e_i,e_j]=C^k_{ij}e_k$. Give $\mathfrak l^*$ the coadjoint action
$(x\lambda)(y)=-\lambda([x,y])$ and form
$\mathfrak l\ltimes\mathfrak l^*[1]$. We use cohomological shifts,
so $[1]$ lowers degree by one. 

Its polynomial CE cochains have
ghosts $c^i$ of degree one and cotangent coordinates $u_i$ of
degree two. The differential is
\begin{equation}\label{mht:eq:centre-ce-signs}
 dc^k=-\frac12C^k_{ij}c^ic^j,
 \qquad du_k=C^\ell_{jk}u_\ell c^j.
\end{equation}
The order $jk$ is the adjoint action on the coefficient $e_k$.
The differential acts from the left by the graded Leibniz rule.

\begin{proposition}\label{mht:prop:finite-centre-coordinates}
Subtract twice the polynomial degree in the $u_i$ from the CE
degree. The substitutions
\[
 u_i\longmapsto O_i,\qquad c^i\longmapsto\partial^i
\]
identify the resulting dg algebra with
\[
 \left(\mathrm{PV}(\operatorname{Sym}\mathfrak l),[\pi,-]_S\right),
 \qquad
 \pi=\frac12C^k_{ij}O_k\partial^i\partial^j.
\]
They preserve the degree-minus-one brackets normalized by
$\{c^i,u_j\}=\delta^i_j$. If $Y_x=\{O_x,-\}$ is the
Hamiltonian derivation for the Kirillov bracket, then
$d_\pi O_x=-Y_x$.
\end{proposition}

\begin{proof}
The regraded coordinates have degrees $|u_i|=0$ and $|c^i|=1$.
Direct Schouten differentiation gives
\[
 d_\pi\partial^k=-\frac12C^k_{ij}\partial^i\partial^j,
 \qquad d_\pi O_k=C^\ell_{jk}O_\ell\partial^j.
\]
These agree with~\eqref{mht:eq:centre-ce-signs}. Both differentials
are derivations, so agreement on generators proves the chain
identity. The brackets also agree on generators, and their
biderivation rule proves the bracket identity. Finally
$\{O_k,O_j\}=C^\ell_{kj}O_\ell$, which is the negative of the
coefficient in $d_\pi O_k$.
\end{proof}

For $[e_1,e_2]=e_2$, these equations read
\[
 dc^1=0,\quad dc^2=-c^1c^2,\quad
 du_1=-u_2c^2,\quad du_2=u_2c^1.
\]
They square to zero. Reversing only the lower indices in the last
two equations instead gives $d^2u_1=-2u_2c^1c^2$. Thus the sign
cannot be changed while retaining the displayed ghost differential.
The regrading is also essential: it does not identify the original
CE degrees with the polyvector degrees of the field theory.

For a formal quantization, fix a filtered continuous
$L_\infty$ quasi-isomorphism $\mathcal U$ from shifted polyvectors
to shifted Hochschild cochains, with first term HKR. This is a
formality datum. If its Taylor series converge in the
$\hbar$-adic filtration, the Poisson tensor defines
\[
 \mu_\hbar=\sum_{r\geq1}\frac{\hbar^r}{r!}
                   \mathcal U_r(\pi,\ldots,\pi),
 \qquad m_\hbar=m+\mu_\hbar.
\]
The Maurer--Cartan equation for $\mu_\hbar$ is associativity of
$m_\hbar$. 

Twisting gives Taylor maps
\[
 \mathcal U^{\hbar\pi}_n(\gamma_1,\ldots,\gamma_n)
 =\sum_{r\geq0}\frac{\hbar^r}{r!}
   \mathcal U_{n+r}(\pi^{\otimes r},\gamma_1,\ldots,\gamma_n).
\]
Their source differential is $\hbar d_\pi$ and their target
differential is $[m_\hbar,-]_G$. The $L_\infty$ equations follow
by collecting coefficients in the original equations after the
substitution $\hbar\pi+\gamma$.

At each parameter quotient the filtration is finite. Its graded
comparison is HKR. Induction on the quotient therefore gives a
quasi-isomorphism. Derived $\hbar$-complete passage gives the
completed comparison. A brace refinement, a PBW identification
with a particular envelope, and compatibility with another
coordinate tower are additional data. In particular the formula
does not construct formality in an arbitrary topological category.

\section{Continuous cochains and the relative center}

For the pointed Hamiltonian algebra of
Chapter~\ref{mht:ch:hamiltonian-coefficients}, put
$R=k[[\hbar]]$, $R_N=R/(\hbar^{N+1})$, and
$A_{d,N}=U_\hbar(\mathfrak h_d)\otimes_RR_N$.
There is no completion in PBW degree at a fixed pair $(d,N)$.
The complete envelope is $\mathcal A=\varprojlim_{d,N}A_{d,N}$.
Its continuous cochains are
\[
 C^n_{\mathrm{uc}}(\mathcal A/R,\mathcal A)
 =\varprojlim_{d,N}\varinjlim_{e\geq d}
       \operatorname{Hom}_{R_N}(A_{e,N}^{\otimes n},A_{d,N}).
\]
The inner index measures input precision and the outer pair
measures output precision. A scalar-valued map may detect a high
Hamiltonian jet, so the input cannot be restricted to the output
cutoff.

Give complete $\mathcal A$-bimodules the exact structure in which
admissible short exact sequences split continuously over $R$.
Proposition~\ref{mht:prop:continuous-relative-bar} constructs the
relative projective bar resolution and proves
\[
 R\operatorname{Hom}^{\mathrm{rel,cont}}_{\mathcal A^e}
                (\mathcal A,\mathcal A)
 \simeq C^\bullet_{\mathrm{uc}}(\mathcal A/R,\mathcal A).
\]
This is a continuous relative center in the specified exact
category. It is not defined by silently changing the tensor
product or the class of admissible exact sequences.

The braces exist because every finite insertion diagram has a
common input refinement. At output $(d,N)$, first choose the
input cutoff of the outer cochain, then use that cutoff as the
output of each inner cochain. A common refinement of their
inputs defines the composite. The usual insertion identities
hold at this refinement and hence at every output.

Antisymmetrization compares the relative cochains with
\[
 \varprojlim_{d,N}\varinjlim_{e\geq d}
 \left(C^\bullet_{\mathrm{CE}}
   (\mathfrak h_e,\operatorname{Sym}\mathfrak h_d)\otimes R_N,
                          \hbar d_{\mathrm{CE}}\right),
\]
where $\mathfrak h_e$ acts through $\mathfrak h_d$. The complete
proof and the derived-limit assertion are
Theorem~\ref{mht:thm:continuous-rees}. This comparison of complexes
does not supply a brace-preserving inverse.

The factor $\hbar$ retains information. If $Z^n$ and $B^n$ are
cycles and boundaries of a vector-space complex $(P,d)$, then
\[
 H^n(P[[\hbar]],\hbar d)=Z^n[[\hbar]]/\hbar B^n[[\hbar]].
\]
In particular its $\hbar$-torsion is the copy of $B^n$ in constant
order. Dividing the differential by $\hbar$ changes this
cohomology before parameter inversion.

The wheel correction has an independent limiting issue. At a
finite Hamiltonian cutoff, the logarithm of the Duflo
half-determinant contains
\[
 \sum_{m\geq1}\frac{B_{2m}}{4m(2m)!}\hbar^{2m}
       \operatorname{Tr}_{\mathfrak h_d}(\operatorname{ad}_x^{2m}).
\]
The traces need not stabilize with $d$. 

The finite-part
subtraction of Proposition~\ref{mht:prop:hamiltonian-even-traces}
gives, on the Cartan line $x_2=aH$,
\[
 \exp(\ell_{\mathrm{ren}}(aH))
          =\frac{\hbar a/2}{\sinh(\hbar a/2)}.
\]
A shift of the cutoff changes the two-wheel finite part. This
defines a subtraction choice. It neither constructs the
continuous formality maps nor proves that their braces agree.

\section{Evaluation at a boundary condition}

Let $\mathcal C$ be a small dg category over $k$, or a specified
compact subcategory of a $k$-linear presentable stable category. Its
diagonal bimodule is
$\Delta_{\mathcal C}(x,y)=\operatorname{Hom}_{\mathcal C}(y,x)$.
Define
\[
 Z_{\mathrm{bimod}}(\mathcal C)
 =R\operatorname{Hom}_{\mathcal C\otimes\mathcal C^{\mathrm{op}}}
                    (\Delta_{\mathcal C},\Delta_{\mathcal C}).
\]
For $b\in\mathcal C$, put $A_b=\operatorname{End}_{\mathcal C}(b)$.
Restriction to the full one-object subcategory gives
\[
 \operatorname{ev}_b:Z_{\mathrm{bimod}}(\mathcal C)
                     \longrightarrow C^\bullet(A_b,A_b).
\]
At the level of bar cochains this simply restricts every object
argument to $b$, so it is a map of cochain complexes and braces.

\begin{proposition}\label{mht:prop:centre-morita-sufficient}
Suppose restriction induces an equivalence of the chosen derived
bimodule categories and carries $\Delta_{\mathcal C}$ to $A_b$.
Then $\operatorname{ev}_b$ is a quasi-isomorphism. In ordinary
derived module categories this holds when $b$ is a compact Morita
generator. Generation is not necessary for the evaluation map
on centers to be a quasi-isomorphism.
\end{proposition}

\begin{proof}
An equivalence preserves derived mapping complexes, which proves
the first assertion. For a compact generator, the tensor--Hom
adjunction with $b$ gives the module equivalence. Its unit and
counit are equivalences on the generator, hence on the localizing
category it generates. Applying the equivalence on both module
variables gives the diagonal bimodule equivalence. The final
assertion follows from the next example.
\end{proof}

\begin{example}\label{mht:ex:centre-without-density}
Let $A$ be the upper triangular $2$ by $2$ matrix algebra. Write
$e_1=E_{11}$, $a=E_{12}$ and $e_2=E_{22}$. For right modules take
$b=e_1A$. Its endomorphism algebra is $e_1Ae_1=k$. The diagonal
has the projective bimodule resolution
\[
 0\longrightarrow Ae_1\otimes e_2A
 \xrightarrow{d}
 (Ae_1\otimes e_1A)\oplus(Ae_2\otimes e_2A)
 \xrightarrow{\mu}A\longrightarrow0,
\]
where
$d(e_1\otimes e_2)=e_1\otimes a-a\otimes e_2$ and $\mu$ is
multiplication. 

The dimensions are $1,4,3$. Multiplication is
onto and its kernel is the displayed one-dimensional image.
Each tensor summand is a direct summand of a free $A^e$-module.

Bimodule Hom into $A$ gives
\[
 k^2\xrightarrow{(\lambda_1,\lambda_2)\mapsto\lambda_1-\lambda_2}k.
\]
Restriction to $b$ is $(\lambda_1,\lambda_2)\mapsto\lambda_1$ in
degree zero and zero in degree one. Its kernel is
$k\xrightarrow{-1}k$, so the restriction is a quasi-isomorphism
to $C^\bullet(k,k)=k$.

Nevertheless $e_2A\ne0$ and
$R\operatorname{Hom}_A(e_1A,e_2A)=e_2Ae_1=0$.
Since $e_1A$ is projective, this vanishing holds in every degree.
It prevents $b$ from generating the derived module category.
Thus equality of center endomorphisms detects less than Morita
generation.
\end{example}

The density quotient $\mathcal C/\langle b\rangle$ and the
cochain cone $\operatorname{Cone}(\operatorname{ev}_b)$ answer
different questions. The former measures generation. The latter
is acyclic exactly when this particular center map is a
quasi-isomorphism. The example has an acyclic cone and a nonzero
density quotient.

For a mixed boundary, one must first construct the category of
boundary conditions on its brane-link site, its complete
factorization bimodules and its diagonal. A restriction functor
must preserve the declared topology and descent conditions.
These are the hypotheses under which the same endomorphism
argument applies. A scalar trace at one Dirac brane does not
provide that bimodule equivalence.

A junction may carry both boundary actions without generating
either boundary category. Section~\ref{mht:sec:ramond-junction}
gives an explicit example. It also distinguishes strict
intertwiners from local operator fields: they have different
derived comparison problems.

\section{Compact support and the direction of a map}

For this section take $k=\mathbb R$ or $\mathbb C$.
Let $\mathfrak l$ be an ordinary Lie algebra and let $I$ be an
open interval. The complex
\[
 \mathfrak l_c(I)=\Omega_c^\bullet(I)\otimes_k\mathfrak l,
 \qquad[\alpha\otimes x,\beta\otimes y]
                   =(\alpha\wedge\beta)\otimes[x,y]
\]
is a dg Lie algebra. Extension by zero gives a covariant map
for inclusions of intervals. For disjoint $I_1,I_2\subset V$,
the two images commute, since their supports are disjoint.
Hence their direct sum maps to $\mathfrak l_c(V)$ as a dg Lie
algebra.

The algebraic CE chains are
$C_*^{\mathrm{CE}}(\mathfrak l_c(I))
=\operatorname{Sym}(\mathfrak l_c(I)[1])$, with direct sums in
symmetric length and with their internal and bracket
differentials. They are covariant. The direct-sum identity for
CE chains therefore gives structure maps
\[
 C_*^{\mathrm{CE}}(\mathfrak l_c(I_1))\otimes
 C_*^{\mathrm{CE}}(\mathfrak l_c(I_2))
 \longrightarrow C_*^{\mathrm{CE}}(\mathfrak l_c(V)).
\]
Their associativity and units follow from those of extension by
zero and direct sum. They define a unital prefactorization
algebra, with $k$ at the empty configuration.

Dual CE cochains run in the other direction. With the algebraic
dual, their structure map is
\[
 C^{\bullet}_{\mathrm{CE}}(\mathfrak l_c(V))
 \longrightarrow
 \operatorname{Hom}_k\left(
 C_*^{\mathrm{CE}}(\mathfrak l_c(I_1))\otimes
 C_*^{\mathrm{CE}}(\mathfrak l_c(I_2)),k\right).
\]
Writing the target as a tensor product of cochain spaces
requires an appropriate dual tensor or a finite-dimensional
hypothesis. With continuous duals the topology must also be
specified. Reversing this arrow is not a factorization
operation.

There is a concrete contraction on compactly supported forms.
Choose $\rho_I\in C_c^\infty(I)$ with integral one. Let $p_I$
integrate a one-form and kill a zero-form, and let
$i_I(1)=\rho_I(t)\,dt$. The target is $k[-1]$, concentrated in
cohomological degree one. Define
\[
 K_I(a(t)\,dt)(t)
 =\int_{-\infty}^t\left(a(s)-\rho_I(s)\int_Ia\right)ds,
 \qquad K_I(f)=0.
\]

The integral in parentheses is zero, so $K_I$ has compact
support. Differentiation proves
\[
 dK_I+K_Id=1-i_Ip_I.
\]
Integration is natural for extension by zero. Thus every
interval inclusion induces a quasi-isomorphism on
$\mathfrak l_c$.

It also induces a quasi-isomorphism on algebraic CE chains.
Filter those chains by length at most $n$. The bracket lowers
length, and the associated quotients are symmetric powers of
the underlying complex. Tensor powers and symmetric
coinvariants preserve quasi-isomorphisms over $k$. Induction
on $n$ proves the assertion on each finite length filtration. Exact filtered colimit proves it for the direct sum. 

This is
a locally constant prefactorization construction. Its
completion and any required Weiss descent must be proved in
the corresponding completed category.

The contraction is not an identification of all current
duals. In the usual test-function topology,
$C_c^\infty(I)'=\mathcal D'(I)$ consists of all distributions
on $I$, while $C^\infty(I)'=\mathcal E'(I)$ consists of
compactly supported distributions. For example integration
against the constant density has noncompact support on $I$
and belongs to the first dual. A cotangent pairing that uses
compactly supported test fields must retain this distinction.

Proposition~\ref{mht:prop:finite-centre-coordinates} can be
applied to an actual finite coefficient model. Extending it
to intervals requires compatible duals, completed tensor
products and transition maps. The equation
$\Phi_V=m(\Phi_{I_1}\otimes\Phi_{I_2})\Delta$ would force
$\Phi_V$ to vanish on $\ker\Delta$. No local coordinate
identity supplies that vanishing. A covariant bulk action
therefore requires its own supported construction.

\section{Cyclic chains and stable matrix traces}

Matrix traces also have a direction. Let $A$ first be an ordinary
unital algebra. Put
\[
 C^\lambda_n(A)=A^{\otimes(n+1)}/(1-t),\qquad
 t(a_0,\ldots,a_n)=(-1)^n(a_n,a_0,\ldots,a_{n-1}).
\]
The cyclic boundary is
\[
 b(a_0,\ldots,a_n)
 =\sum_{i=0}^{n-1}(-1)^i
    (a_0,\ldots,a_ia_{i+1},\ldots,a_n)
   +(-1)^n(a_na_0,a_1,\ldots,a_{n-1}).
\]
Over $k$, this cyclic coinvariant complex computes cyclic
homology. 

For a dg algebra the rotation includes the internal
Koszul sign, and the internal differential is totalized with $b$.
The reduced theory of an augmented algebra is the kernel of the
map to the corresponding complex of $k$.

Use the homological CE convention
$d(X\wedge Y)=[X,Y]$. The alternating matrix trace is
\begin{equation}\label{mht:eq:matrix-cyclic-trace}
 \begin{split}
 \gamma_{N,n}(X_1\wedge\cdots\wedge X_n)
   =\sum_{\sigma\in S_{n-1}}\operatorname{sgn}(\sigma)
     \sum_{i_1,\ldots,i_n}
     (&(X_1)_{i_1i_2},
       (X_{\sigma(2)})_{i_2i_3},\ldots,\\[-1mm]
      &(X_{\sigma(n)})_{i_ni_1})
       \quad\text{in }C^\lambda_{n-1}(A).
 \end{split}
\end{equation}
Here $S_{n-1}$ permutes $2,\ldots,n$. For $n=1$ the formula is
the ordinary matrix trace. It is equivalent to full
antisymmetrization divided by $n$, using the cyclic relation.

\begin{proposition}\label{mht:prop:matrix-trace-chain}
For every $N$, the maps~\eqref{mht:eq:matrix-cyclic-trace} satisfy
$b\gamma_N=\gamma_Nd$. They are compatible with stabilization by
a zero block. The same statement holds for dg algebras with the
total Koszul signs.
\end{proposition}

\begin{proof}
The cyclic relation makes the expression alternating in all
$n$ inputs. A face of its boundary multiplies two successive
matrix entries, and summing the intermediate index is matrix
multiplication. For each unordered pair of inputs, the two
orders give their commutator. The permutation signs are the
signs that move this pair to the front of the CE wedge. The
remaining permutations are exactly those in the alternating
trace of that commutator and the remaining inputs. 

The closing
cyclic face supplies the pair across the chosen starting point.
This partitions every summand on both sides and proves the
identity. Zero blocks add only zero entries. 

For homogeneous
dg inputs, each reordering moves the same suspended inputs on
both sides. The internal differential commutes with products
and trace. This proves the total version.
\end{proof}

The reverse sum of cyclic matrix tensors has a different
normalization. Define
\[
 \Theta_1(a)=\sum_i E_{ii}a,\qquad
 \Theta_2(a,b)=\sum_{i,j}E_{ij}a\wedge E_{ji}b.
\]
Then
\begin{equation}\label{mht:eq:raw-cyclic-rank-defect}
 d\Theta_2(a,b)=N\Theta_1(ab-ba).
\end{equation}

Indeed the forward product sums to $N\sum_iE_{ii}ab$ and the
reverse product sums to $N\sum_jE_{jj}ba$. Thus the unnormalized
$\Theta$ is not a chain map from the cyclic complex for a
noncommutative $A$ at $N>1$. The factor is already visible for
two free generators. A finite rank bound does not remove it.

There is a chain-level primitive projection after taking the
correct stable coinvariants. Let
\[
 \mathfrak{gl}_\infty(A)=\varinjlim_N\mathfrak{gl}_N(A),\qquad
 H(A)=C_*^{\mathrm{CE}}(\mathfrak{gl}_\infty(A))_{
                                      \mathfrak{gl}_\infty(k)}.
\]
The subscript means coinvariants for the scalar matrix Lie
algebra. Stabilization and these coinvariants make block sum
an associative, graded-commutative chain product. The CE
coproduct makes $H(A)$ a connected dg Hopf algebra.

\begin{theorem}[Stable cyclic primitives]
\label{mht:thm:stable-cyclic-primitives}
The primitive subcomplex of $H(A)$ is naturally isomorphic to
$C^\lambda_{*-1}(A)$. The Hopf logarithm
\[
 P_{\mathrm{prim}}=\log^*(1)
 =\sum_{r\geq1}\frac{(-1)^{r-1}}r
                    (1-u\varepsilon)^{*r}
\]
is a chain idempotent onto that subcomplex. Its sum is finite
in each CE length. Consequently
\[
 \operatorname{Prim}H_*^{\mathrm{CE}}
                  (\mathfrak{gl}_\infty(A))\cong HC_{*-1}(A).
\]
For an augmented algebra, put
\[
 H_{\mathrm{rel}}(A)=H(A)/(H(k)^+H(A)),
\]
where $H(k)^+$ is the augmentation ideal of the scalar Hopf
subalgebra. Its primitive complex is
$C^\lambda_{*-1}(A)/C^\lambda_{*-1}(k)$, giving reduced cyclic
homology.
\end{theorem}

\begin{proof}
At CE length $n$, scalar matrix coinvariants have the
permutation description
\[
 \bigl(k[S_n]\otimes A^{\otimes n}\otimes\operatorname{sgn}
                                           \bigr)_{S_n}.
\]
The symmetric group conjugates the permutation label and
permutes the $A$-inputs. This follows from matrix invariant
theory: contractions are indexed by permutations, and they
are independent at matrix rank at least $n$.

The coproduct splits the labels into subsets preserved by
the permutation. A one-cycle permutation is therefore
primitive. Every permutation is a disjoint collection of
cycles, so the Hopf algebra is the graded symmetric algebra
on the one-cycle subspace. Identifying all starting points
of one cycle gives exactly the cyclic relation defining
$C^\lambda_{n-1}(A)$.

For a cycle, choose the class represented by
$E_{12}a_1\wedge E_{23}a_2\wedge\cdots\wedge E_{n1}a_n$,
using distinct matrix indices. Its alternating trace is
$(a_1,\ldots,a_n)$. The CE boundary joins adjacent entries,
including the last and first, with the signs of $b$.
It follows that the primitive identification is an
isomorphism of complexes. This is the chain calculation in
the Loday--Quillen theorem.

On a product of $n$ positive primitives, the Hopf logarithm
has coefficient
\[
 \sum_{r=1}^n(-1)^{r-1}(r-1)!S(n,r),
\]
where $S(n,r)$ counts partitions into $r$ nonempty parts.
The coefficient is $1$ for $n=1$ and $0$ otherwise, as follows
from $\log(1+(e^t-1))=t$. This proves idempotence and its
image. All the maps in its convolution formula are chain
maps.

Finally the scalar matrix action on CE chains is nullhomotopic
by insertion of a scalar matrix. In each finite rank its
rational representations are semisimple. The nontrivial
isotypic summands have zero homology, so passage to
coinvariants is a quasi-isomorphism. Filtered colimit gives
the stable assertion. 

The algebra augmentation splits the
scalar cyclic subcomplex. The symmetric-algebra description
then identifies the Hopf quotient by $H(k)^+$ with the symmetric
algebra on the reduced cyclic complex. This proves the relative
statement.
\end{proof}

For dg $A$ concentrated in nonnegative homological degrees,
the proof uses total Koszul signs and the same first-quadrant
filtrations. At a fixed CE length $n$ the scalar matrix
coinvariant description is valid for $N\geq n$. Computing
homology through length $L$ also uses the next chain term.
Thus $N\geq L+2$ is a sufficient bound. 

A separate polynomial
trace window of weight $K$ can impose $N\geq K$. The common
bound $N\geq\max(K,L+2)$ is sufficient for both purposes.
The projection acts on the specified coinvariant complex.
It is not deletion of multitrace tensors in the original
fixed-rank CE complex.

\section{The full cyclic-word resolution}

The scalar quotient in cyclic homology must also be
distinguished from removing an antifield class. Use
homological degrees in this section and set
\[
 R=k\langle x,y,\psi\rangle,\qquad
 |x|=|y|=0,\quad |\psi|=1,\quad
 d\psi=xy-yx.
\]
The free-resolution contraction of
Proposition~\ref{mht:prop:word-contraction} identifies $R$ with
$k[x,y]$ in the derived category of associative algebras.
Define its reduced cyclic-word complex by
\[
 R_{\mathrm{cyc,red}}=R/(k1+[R,R]_{\mathrm{gr}}).
\]
The denominator is a vector subspace of graded commutators,
not a two-sided ideal. 

Give $x,y,\psi$ bidegrees $(1,0)$,
$(0,1)$ and $(1,1)$, respectively. The differential preserves
this bidegree.

\begin{theorem}\label{mht:thm:full-cyclic-word-homology}
The cyclic-word homology is
\[
 H_p(R_{\mathrm{cyc,red}})=
 \begin{cases}
 k[x,y]/k,&p=0,\\
 k[x,y]\,dx\wedge dy,&p=1,\\
 0,&p\geq2.
 \end{cases}
\]
Let $\mathsf W_{a,b}$ consist of words with $a$ copies of $x$
and $b$ copies of $y$, for $a,b\geq0$.
In bidegree $(a+1,b+1)$, the degree-one group is generated by
\[
 \Psi_{a,b}=\sum_{w\in\mathsf W_{a,b}}[\psi w].
\]
In particular $[\psi]$ survives.
\end{theorem}

\begin{proof}
First compare cyclic words with cyclic chains of the free
dg algebra $R$. For a graded vector space $V$, the tensor
algebra $T(V)$ has the bimodule resolution
\[
 0\longrightarrow T(V)\otimes V\otimes T(V)
 \longrightarrow T(V)\otimes T(V)
 \longrightarrow T(V)\longrightarrow0.
\]
The first map sends $a\otimes v\otimes b$ to
$av\otimes b-a\otimes vb$. At word length $m>0$, its
Hochschild kernel and cokernel are the invariants and
coinvariants of the graded cyclic rotation on $V^{\otimes m}$.

The Connes operator between them is the norm
$1+t+\cdots+t^{m-1}$. In characteristic zero that norm is an
isomorphism from coinvariants to invariants. Thus the reduced
cyclic complex of a tensor algebra has only its cyclic-word
quotient in cyclic degree zero.

Apply this calculation to the underlying graded tensor
algebra of $R$, and then restore its internal differential.
The cyclic double complex is first quadrant in homological
degrees. In each positive bidegree there are only finitely
many words and internal degrees. The comparison therefore
gives
\[
 H_*(R_{\mathrm{cyc,red}})=\overline{HC}_*(R).
\]
Tensor powers preserve the quasi-isomorphism
$R\to k[x,y]$ over $k$. The same first-quadrant comparison
identifies this with $\overline{HC}_*(k[x,y])$.

For $k[x,y]$, the diagonal Koszul resolution gives
Hochschild chains equal in homology to polynomial forms.
The normalized map
$a_0[a_1|\cdots|a_n]\mapsto a_0da_1\wedge\cdots\wedge da_n/n!$
intertwines the Connes operator with the de Rham differential.
The cyclic bicomplex consequently gives
\[
 \overline{HC}_0=k[x,y]/k,\qquad
 \overline{HC}_1=\Omega^1/dk[x,y],
\]
and zero in higher degrees. To see the last assertion
directly, contract positive-weight polynomial forms by the
Euler vector field divided by total weight. The only
uncontracted de Rham class is the scalar class, removed in
the reduced complex. Differentiation identifies
$\Omega^1/dk[x,y]$ with $\Omega^2=k[x,y]dx\wedge dy$.

It remains to identify the cycle. Replacing the displayed
$\psi$ by $xy-yx$ shows that $d\Psi_{a,b}=0$: every
zero-antifield cyclic word has equally many cyclic $xy$ and
$yx$ transitions. Define a functional on one-antifield
words by $[\psi w]\mapsto1$ in this fixed bidegree. A
two-antifield boundary has the four terms
\[
 d[\psi u\psi v]
   =[\psi(vxyu-vyxu-uxyv+uyxv)],
\]
whose functional values sum to zero. 

The value on
$\Psi_{a,b}$ is $\binom{a+b}{a}\ne0$. Hence its class is
nonzero and spans the already computed one-dimensional
group. This also specifies its normalization: the
functional divided by $\binom{a+b}{a}$ sends $\Psi_{a,b}$
to $1$.
\end{proof}

The class $[\psi]$ has bidegree $(1,1)$ and corresponds to
the constant area form. It is not the scalar augmentation
class of bidegree $(0,0)$. Removing it is an additional
quotient of the cyclic-word complex. In the full cyclic
chain complex, its representative includes the bar term
that cancels $d\psi=xy-yx$. With the convention
$b(x|y)=xy-yx$, it is $\psi-(x|y)$. Quotienting only the
two-term subcomplex $\psi\mapsto xy-yx$ does not remove
this cyclic homology class.

\section{Primitive projections and the scalar rank}

Let $L^\nu$ be the vector space of cyclic words in
$k\langle x,y\rangle$, including the empty cyclic word
$\nu$. Define its necklace Lie bracket by cutting an $x$ in the first
word and a $y$ in the second. Join the remaining words and
subtract the reverse cuts. The empty word is
central and
\[
 [x,y]=\nu.
\]

Associativity of concatenation and pairing opposite orders
of two cuts prove Jacobi. The symmetric algebra
$\operatorname{Sym}(L^\nu)$ has the induced linear Poisson
bracket. Its Hopf coproduct makes every element of $L^\nu$
primitive, including $\nu$.

Evaluation on matrices sends $\nu$ to $\operatorname{Tr}I_N=N$.
Before this specialization, $\nu$ counts one trace. After it,
the same term is a scalar. This change of trace count affects
the primitive projection.

\begin{proposition}\label{mht:prop:primitive-rank-bracket}
Projection $\operatorname{Sym}(L^\nu)\to L^\nu$ onto symmetric
degree one is a Lie map. The quotient $L=L^\nu/k\nu$ also
gives a Lie projection $\operatorname{Sym}(L)\to L$.
Abelianization of the nonempty cyclic words then gives a
Lie map $L\to k[x,y]/k$.

At nonzero fixed rank, the projection that kills scalar and
multitrace terms is not a Lie map for the unreduced matrix
Poisson bracket. On $X=\operatorname{Tr}x$ and
$Y=\operatorname{Tr}y$, its defect is
\begin{equation}\label{mht:eq:primitive-rank-defect}
 \pi_{\mathrm{prim}}\{X^2,Y\}
 -\{\pi_{\mathrm{prim}}X^2,\pi_{\mathrm{prim}}Y\}
                         =2N X.
\end{equation}
\end{proposition}

\begin{proof}
The linear Poisson bracket sends symmetric degrees $(r,s)$
to degree $r+s-1$. With positive degrees, its degree-one
part occurs only at $(1,1)$. Scalars have zero bracket.
This proves the first projection identity. The central Lie
quotient preserves it. On cyclic words, abelianization turns
the cut formula into the ordinary Hamiltonian bracket,
modulo its constant term.

At fixed rank, $\{X,Y\}=N$ and the Leibniz rule gives
$\{X^2,Y\}=2NX$. The projection kills $X^2$, retains $X$
and $Y$, and kills the scalar bracket of $X$ and $Y$.
This proves~\eqref{mht:eq:primitive-rank-defect}.
\end{proof}

For $N\geq2$, the degree-two trace window already separates
the single-trace and double-trace terms in this example.
The class $X$ is nonzero on the commuting-matrix locus.
Thus the defect cannot be made into a Koszul boundary by a
homotopy between the same two projected bracket maps. The
zero-central-character construction precedes specialization. It does not quotient the fixed-rank scalar $N1$ to zero in
a nonzero unital algebra.

The Hamiltonian abelianization extends to a chain map from
the full cyclic-word complex by killing words containing
$\psi$. A one-antifield differential becomes a commutator,
which abelianizes to zero. Higher-antifield terms still
contain an antifield after one differential. Extending over
the symmetric algebra and projecting to its linear part
therefore preserves the differential. Its bracket scope
is the one just proved. The independent one-antifield
classes of Theorem~\ref{mht:thm:full-cyclic-word-homology}
must be retained for a cotangent comparison.

There is a useful product truncation as well. For any
symmetric algebra $S(L)$, projection to $k\oplus L$ is the
quotient by $S^{\geq2}(L)$, with square-zero multiplication
on $L$. If $F,G\in S(L)$, then
\[
 \pi_1(FG)=\varepsilon(F)\pi_1(G)+\pi_1(F)\varepsilon(G).
\]
This is the linear part of the product with its scalar
unit retained. A Lie bracket on $L$ is a separate
operation. For supported observables, a factorization
interpretation of this quotient also requires compatible
structure maps and descent on the chosen site.

\section{Actions on the ordered collision carrier}

The geometric center problem on a curve has a different
recipient. We specify it over $X=\mathbb A^1_k$. For a finite
set $I$, form a scheme $\mathscr R_I$ from one copy of $X^I$
for each total order of $I$. Two copies are glued where
$x_i\ne x_j$ for every pair whose orders disagree. 

Distinct
positions forget their relative order. Coincident positions
retain it. For two labels there are two diagonal branches.

An ordered surjection $I\twoheadrightarrow J$ includes an
order on each fiber. Replacing each label of an order on $J$
by its ordered fiber gives a collision map
$\mathscr R_J\to\mathscr R_I$. Let $\mathcal D$ be the
category of compatible derived right differential modules
on these schemes, with extraordinary pullback along every
such collision map. Compatibility includes its composition
homotopies. It retains full normal direct images, including
normal derivatives.

There are two tensor operations. Disjointness convolution $\star$ sums over ordered
decompositions $I=I_1\sqcup I_2$. For each decomposition,
externally tensor the coefficients on the disjoint open
and apply open direct image.
Its unit $K$ is supported at the empty configuration. 

The
diagonal tensor is
\[
 (F\otimes^!G)_I=\Delta^!(F_I\boxtimes G_I).
\]
Its unit $J$ is the compatible dualizing family. On a chart
over $X^I$, $J_I=\omega_{X^I}[|I|]$. The map $K\to J$ is
not an identification of the two units.

Use connected coaugmented counital coassociative coalgebras
for $\star$, with empty coefficient $k$. Denote their
category by $\mathcal C$. The forward interchange keeps
common ordered decompositions in
\[
 (F\star G)\otimes^!(H\star L)
       \longrightarrow(F\otimes^!H)\star(G\otimes^!L).
\]
It tensors coalgebra coproducts. Consequently
$(\mathcal C,\otimes^!,J)$ is symmetric monoidal.

A coalgebra is \emph{factorizing} if its coproduct becomes
an equivalence for each ordered decomposition on its
disjoint open. Write $\mathcal C^{\mathrm f}$ for this
full subcategory.

A transverse monoid is a unital associative algebra object
in $(\mathcal C,\otimes^!,J)$. Its product and unit are
coalgebra maps. An action on $C$ consists of
$\alpha:F\otimes^!C\to C$ and all the unit and
associativity homotopies of a module structure. Morphisms
of actions preserve these comparisons. A \emph{factorizing
monoid center} of $C\in\mathcal C^{\mathrm f}$ means a
final object among these actions with $F\in\mathcal C^{\mathrm f}$.

One can first allow all coalgebras as actors. The resulting
endomorphism object always exists on this carrier. 

Positive
objects have zero empty coefficient, so their $r$-fold
convolution vanishes on configurations of fewer than $r$
positions. The normalized associative bar sums suspended
words in the augmentation ideal. Its differential multiplies
adjacent factors and its coproduct cuts words. The normalized
cobar sums desuspended words in the reduced coalgebra, with
differential induced by its coproduct. Both sums are finite
at each configuration.

To compare their composites with the original object,
filter by the number $n$ of original factors. Each of the
$n-1$ gaps is either inside a block or between blocks.
On associated graded its two states form an identity
complex, with the suspension signs. Their tensor product
is contractible for $n>1$. Lengths zero and one retain the
unit and original reduced object. 

The filtration ends at
$n=|I|$ on a configuration $I$. Thus both comparisons are
equivalences. Finite normalization commutes with collision
pullback, so the comparisons and their coherences belong
to the declared ordered carrier.

The category of augmented convolution algebras is
presentable, since nonunital algebras on positive objects
are algebras for an accessible tensor-algebra monad.
Bar--cobar gives presentability of $\mathcal C$. Coalgebra
colimits are computed on the underlying reduced families:
equivalently, they are limits of nonunital algebras in the
opposite category. The diagonal tensor preserves these
colimits. The adjoint functor theorem therefore supplies
an internal Hom $[C,D]_!$ in $\mathcal C$.

\begin{proposition}\label{mht:prop:ambient-actor}
The object $\mathcal E(C)=[C,C]_!$ has a coherent
composition product and unit $J\to\mathcal E(C)$.
Its evaluation is a transverse action, and
\[
 \operatorname{Map}_{\operatorname{Alg}(\mathcal C)}
                        (F,\mathcal E(C))
 \simeq\{\text{coherent transverse actions of }F\text{ on }C\}.
\]
If $\mathcal E(C)$ factorizes, it is a factorizing monoid
center. Factorization of this internal Hom is an additional
condition.
\end{proposition}

\begin{proof}
Adjoint to two successive evaluations is a product
$\mathcal E(C)\otimes^!\mathcal E(C)\to\mathcal E(C)$.
The identity of $C$ gives its unit. Iterated evaluation is
independent of parentheses with all its coherent
comparisons. Adjunction carries the multiplication and
unit equations of a map into $\mathcal E(C)$ to the action
equations, and carries their higher simplices in the same
way. This proves the mapping-space equivalence. If the
object factorizes, this universal action belongs to the
restricted category as well and remains final there.
\end{proof}

A pointed actor carries only a map $u:J\to F$. Its pointed
action space is
\[
 \Phi_C(F,u)=\operatorname{Map}_{C/}(F\otimes^!C,C),
\]
where the structure map from $C$ is $u\otimes^!1$.
Representing this functor on factorizing pointed actors is
sufficient to construct a factorizing monoid center by
composition of the representing action. A pointed action
does not itself include the multiplication homotopies of
a monoid action. Their representability questions must
therefore be compared at their actual morphism scopes.

\section{The center under a cohomological bound}

Let $Q$ be an ordinary unital chiral algebra, flat over
$\mathcal O_X$, and write $C_Q$ for its
Chevalley--Cousin factorization coalgebra on the ordered
carrier. Its main-diagonal Lie object $L_Q$ has singleton
right coefficient $\omega_X\otimes Q$, and
$C_Q=\operatorname{CE}_\star(L_Q)$.

The chiral center $Z(Q)$ consists of the states whose
singular products with every state vanish. It is a
commutative differential algebra. The tensor product of
ordinary chiral algebras represents pairs of commuting
maps: this is the Beilinson--Drinfeld tensor-product
property. In particular a map $B\otimes Q\to Q$ that is
the identity on $Q$ is a map $B\to Z(Q)$.

For a derived actor, the corresponding bound is the
ordinary factorization heart. On every ordered chart
require
\[
 F_I\in D^{\leq-|I|}(\mathcal D_{X^I}).
\]
After the dualizing shift and passage to left modules,
its top cohomology has degree zero. This condition permits
negative derived degrees but excludes positive degrees
above that heart.

\begin{theorem}\label{mht:thm:bounded-geometric-centre}
For pointed factorizing actors satisfying this bound,
restriction of an action on $C_Q$ to its vacuum gives an
equivalence of discrete spaces
\[
 \Phi_{C_Q}(F,u)
     \simeq\operatorname{Map}_{J/}(F,C_{Z(Q)}).
\]
The central action of $C_{Z(Q)}$ is therefore final among
the factorizing transverse monoid actions satisfying
the same bound.
\end{theorem}

\begin{proof}
Maps from an object below the heart into an object of the
heart are determined on top cohomology and have discrete
mapping spaces. The same holds for the compatibility
limits defining coefficient maps and their unit
conditions. Write $M_I$ for the top left module of $F_I$
and $A_I$ for the left factorization module of $Q$.
An action gives $a_I:M_I\otimes A_I\to A_I$.

Let $e$ be the image of the actor vacuum and set
$\zeta(f)=a_1(f\otimes1)$. The pointing says
$a_1(e\otimes q)=q$. The two-position localization
triangle supplies a chiral boundary operation $\mu_F$
on $M_1$. Naturality of its connecting morphism under
$a_2$ gives
\[
 \mu_Q((a_1\boxtimes a_1)(x\otimes c))
                 =a_1(\mu_F(x)\otimes c|_\Delta).
\]
Here $x$ is a meromorphic section of $M_1\boxtimes M_1$
on the disjoint open and $c\in A_2$.

Take $x=(f\boxtimes e)\varphi$ and use the two vacuum
sections $1\boxtimes q$ and $q\boxtimes1$ of $A_2$.
Their restrictions to the diagonal agree. Hence
\[
 \mu_Q((\zeta(f)\boxtimes q)\varphi)
 =\mu_Q((a_1(f\otimes q)\boxtimes1)\varphi).
\]
At $\varphi=1$ the right side is zero in the full normal
direct image. Its normal-derivative coefficients are all
singular products, so $\zeta(f)\in Z(Q)$.
At $\varphi=(z_1-z_2)^{-1}$ the
vacuum residue axiom gives
$a_1(f\otimes q)=\zeta(f)q$.

Restriction to the vacuum thus gives a coalgebra map
$g:F\to C_Q$ whose singleton image lies in $Z(Q)$.
Factorization puts its image in the external powers of
$Z(Q)$ on every disjoint open. The Chevalley--Cousin
kernel description gives
\[
 (C_{Z(Q)})_I
   =(C_Q)_I\cap j_{I,*}j_I^*(Z(Q)^{\boxtimes I})
\]
in the corresponding left-module description. This is
because the first Cousin differential on the central
subalgebra is its own Cousin differential. Thus $g$
factors through $C_{Z(Q)}$. 

The factorization modules
of the flat chiral algebra have no sections supported
only on diagonals. Agreement on the disjoint open
therefore proves that its central action recovers
$\alpha$. Evaluation on the vacuum recovers $g$ in
the opposite direction.

For a monoid actor, associativity of its action now says
that $g$ preserves multiplication. The action unit says
that it preserves the unit. All mapping spaces involved
are discrete, so these equations determine their higher
coherences uniquely. This proves finality in the
bounded category.
\end{proof}

The cohomological bound precedes the use of the residue.
Without it, top cohomology can forget action homotopies.
For example $\pi_1\operatorname{Map}(k[-1],k)=k$,
although $H^0(k[-1])=0$. The following construction
shows that, for ordinary boundaries with a derivation,
removing the bound changes the existence question itself.

\section{A factorizing monoid obstruction}

Work now over $k=\mathbb C$. Let $Q$ be an algebraic
degree-zero unital vertex algebra, and fix a nonzero
vertex derivation $D$. Thus $D1=0$, $DT=TD$, and
$D(x_{(n)}y)=(Dx)_{(n)}y+x_{(n)}Dy$ for every integer
$n$. Put $C=C_Q$ and $\ell=L_k$.

For a finite-dimensional vector space $W$, place its
exterior generators in degree one and define
\[
 E_W=\Lambda(W[-1]),\qquad
 C_W=\operatorname{CE}_\star(\ell\otimes E_W).
\]
The differential and translation on the exterior
generators are zero. These explicit actors have an
augmentation to $J=C_0$. Arbitrary receiving actors below
are not required to have one.

We need natural maps of these exterior actors. A natural
operation uses products, Lie brackets, permutations and
exterior coefficient maps, with each ordinary Lie input
used once. Its identities hold for every test Lie
algebra. Complete only in the number of exterior input
generators. At $m$ positions this number is at most
$m\dim W$, so evaluation remains finite on the original
coefficient diagrams.

\begin{lemma}\label{mht:lem:marked-actor-maps}
Natural maps $C_W\to C_V$ under and over $J$ are
classified, including their natural homotopies, by
formal matrices
\[
 A(t)\in\operatorname{Hom}(W,V)[[W^*]].
\]
Their components are contractible. The matrix remembers
which input supplies the exterior output. With
$\Phi_A(t)=A(t)t$, composition and tensor product are
\[
 (B\circ A)(t)=B(\Phi_A(t))A(t),\qquad
 A\otimes B=\begin{pmatrix}A&0\\0&B\end{pmatrix}.
\]
Thus an associative unital multiplication of these
matrices defines a coherent transverse monoid on $C_W$.
\end{lemma}

\begin{proof}
The scalar-relative convolution algebra of natural
operations has a weightwise model obtained from the
Lie cobar of $(E_W^+)^*$. In a basis its generators
$y_S$, for nonempty exterior subsets $S$, have degree
$1-|S|$. The differential splits $S$ into two subsets
and takes their signed Lie bracket. An extra scalar
generator $z$ has degree one and $dz=z^2$ in the
associative envelope.

The ideal generated by $z$ contracts by
$h(z^{2j})=z^{2j-1}$ and $h(z^{2j-1})=0$, applied
to the first such factor in a free-product word.
The remaining Lie cobar projects onto the abelian
Lie algebra $W^*$. Its envelope is a quasi-isomorphism
onto $\operatorname{Sym}(W^*)$: the dual Koszul
resolution is
\[
 E_W\otimes k[v_1,\ldots,v_r],\qquad dv_i=\alpha_i,
\]
and positive total weight $n$ contracts by
$n^{-1}\sum_i v_i\partial_{\alpha_i}$.
PBW symmetrization identifies the Lie length-one
summands as well.

Tensoring with $E_V$ and taking the positive target
exterior ideal gives a surjective dg algebra map to
\[
 \left[\widehat{\operatorname{Sym}}(W^*\otimes E_V)
                                            \right]_{E_V^+}
\]
with acyclic kernel in every source weight. This
commutative model has no nonpositive degrees. Its
degree-one part is
$\widehat{\operatorname{Sym}}(W^*)\otimes W^*\otimes V$.
The Maurer--Cartan equation is automatic there.
Weightwise lifting through the acyclic kernel, also
with polynomial forms relative to simplex boundaries,
gives the contractible components.

The remaining degree-one symbol has one marked input.
Substitution into a second natural map first replaces
the ordinary coordinate by $A(t)t$, then carries its
marked input by $A(t)$. This proves the composition
formula. Independent exterior factors give the
block-diagonal formula. These formulas identify the
natural category with the nerve of the stated matrix
category. Associativity and unit equations lift with
all their coherences through the contractible fibers.
\end{proof}

Use the auxiliary dual-number algebra with basis $a,b$
and products $a^2=a$, $ab=ba=b$, $b^2=0$. Its unit $a$
is not the vacuum of an actor. Let $L_x$ be multiplication
by $x$ in this two-dimensional algebra and set
\[
 Q(x)=e^{L_x},\qquad
 M(x)=\sum_{r\geq0}\frac{L_x^r}{(r+1)!}.
\]
The constant term of $M$ is the identity. 

The two
blocks
\begin{equation}\label{mht:eq:actor-frame-product}
 \bigl(M(x+y)^{-1}M(x)\mid
        M(x+y)^{-1}Q(x)M(y)\bigr)
\end{equation}
define a multiplication on
$F=C_{ka\oplus kb}$. Its coordinate contraction is
addition. 

To verify associativity, use fiber coordinates
$U=M(x)u$. The product becomes
\[
 (x,U)(y,V)=(x+y,U+Q(x)V).
\]
The identity $Q(x+y)=Q(x)Q(y)$ proves associativity.
The identity $M(x)x=e^x-a$ proves that the diagonal
$u=x$ is closed under this product, giving the claimed
coordinate contraction.

The lines $G=C_{ka}$ and $H=C_{kb}$ are submonoids.
On $H$ the product is additive, since $L_{sb}b=0$.
Write $i:G\to F$ and $j:H\to F$ for their inclusions.

\begin{lemma}\label{mht:lem:relative-derivation-actor}
There is a coherent unital action of $F$ on $C_Q$.
Its restriction to $G$ is trivial, and its restriction
to $H$ has first derivation term $D$.
Both inclusions are morphisms of monoid actions.
\end{lemma}

\begin{proof}
Tensoring with $C_Q$ is strong symmetric monoidal
for $\star$. On an ordered disjoint decomposition
its comparison is the factorization equivalence
of $C_Q$, followed by open projection. Iteration
retains the same ordered partition, proving all
unit and collision comparisons. On main diagonals
it carries $\ell$ to $L_Q$, including its chiral
bracket. Thus it carries $C_W$ to
$\operatorname{CE}_\star(L_Q\otimes E_W)$ and
commutes with the positive bar construction.

The vertex automorphism $1+vD$ over $k[v]/(v^2)$
shows that these comparisons carry $D$ to
$D\otimes1$ after exterior extension. Natural
differential polynomials in $D$ therefore evaluate
on the full collision diagrams.

For $r$ actors let $t_\nu,s_\nu$ be dual to the
$a,b$ directions. Give $t_\nu,s_\nu,D$ auxiliary
weights $0,-1,1$, respectively. Prescribe the
entire action to be trivial when every $b$ input
is absent. The natural relative cochain complex
then has the commutative model
\[
 \left[(s_1,\ldots,s_r)
 k[[t_1,s_1,\ldots,t_r,s_r]]\otimes
                 \Lambda(kD[-1])\right]^{\mathbb G_m}.
\]

Here $D[-1]$ has cohomological degree one. To
obtain this model, express universal differential
polynomials as convolution products of unary
coefficient extractions and divided derivatives.

After the contractible scalar generator is removed,
their differential is the cobar differential of
$\Lambda(P[1])\otimes\Gamma(kD)$, where $P$
records the actor coordinates. The elementary
exterior and polynomial Koszul resolutions give
$\widehat{\operatorname{Sym}}(P)\otimes\Lambda(kD[-1])$.
The projection is surjective with acyclic kernel
weight by weight.

Restriction to the prescribed ideal and auxiliary
weight is exact. All deletions of actor or exterior
coordinates split and commute. The same projection
therefore has acyclic kernel relative to their
compatible prescribed faces. Its target has no
nonpositive degrees, and its degree-one terms are
exactly $\sum_\nu s_\nu f_\nu(t_1,\ldots,t_r)D$.
Each corresponding relative action component is
contractible by weightwise Maurer--Cartan lifting.
The bar and Chevalley versions give the same
lifting complex: their PBW contraction is finite
on each input and preserves every exterior face.

Choose $sD$. Pullback by~\eqref{mht:eq:actor-frame-product}
gives $(s_1+s_2)D$. Composition gives the same
class, because its operator symbol is
$e^{s_1D}e^{s_2D}=e^{(s_1+s_2)D}$.
Relative contractibility fills the action
associahedra and their unit faces. It retains the trivial
$G$ action and the prescribed derivation action of $H$. 

Every coefficient
uses finitely many operations at a configuration.
Evaluation through the strong tensor comparison
proves the claimed geometric action and maps.
\end{proof}

Let $S=C_{k^2}$. Insert its two exterior generators
into four copies of $F$ as
$\varepsilon_1\mapsto a^{(1)}-a^{(3)}$ and
$\varepsilon_2\mapsto b^{(2)}-b^{(4)}$, then multiply
in that order. The resulting pointed map $w:S\to F$
factors, up to pointed homotopy, through $j$ as
$w=jf$. The marked matrix of $f:S\to H$ is
\begin{equation}\label{mht:eq:actor-marked-commutator}
 A_f(t,s)=\bigl(-s\chi(t)\mid t\chi(t)\bigr),
 \qquad \chi(t)=\frac{e^t-1}{t}
                 =\sum_{r\geq0}\frac{t^r}{(r+1)!}.
\end{equation}

Indeed the affine commutator has zero coordinate
and translation $\chi(t)(tq-sr)b$ on fiber inputs
$ra,qb$. Negating coordinate and fiber gives the
inverse since $M(-x)=Q(-x)M(x)$. Lemma~\ref{mht:lem:marked-actor-maps}
then identifies the pointed maps. The contraction
$A_f(t,s)(t,s)^T$ vanishes, but its first marked
form $t\,ds-s\,dt$ is nonzero.

\begin{lemma}\label{mht:lem:geometric-actor-detector}
Let $h:H\to Z$ be a pointed map receiving the
derivation action, with $Z$ any factorizing
coassociative actor on $C_Q$. No cohomological bound,
cocommutativity, or augmentation to $J$ is required
of $Z$. The pointed map $hf$ is not homotopic to
the map $S\to J\to Z$.
\end{lemma}

\begin{proof}
Put $B=\Omega Z$ and let $h_0$ be the twisting
component on one pure exterior actor input. Its
plain singleton class lies in
$H^1R\operatorname{Hom}_X(\ell,B_{\{1\}})$.
We first show that it is nonzero on every nonempty
affine open $Y$.

A degree-zero class $\ell|_Y\to B_{\{1\}}|_Y$
acts on $L_Q|_Y$ by an endomorphism $M$. It obeys
\begin{equation}\label{mht:eq:actor-module-detector}
 M\mu_Q=\mu_Q(1\star M)=\mu_Q(M\star1).
\end{equation}

To see this without a condition on $B$, restrict
to two distinct roots. Both ordered cobar products
are equivalences there by factorization. Subtract
the open exchange commutator and its action
nullhomotopy. The remainder factors through
\[
 \operatorname{cofib}
   (j_!j^*(\ell\boxtimes\ell)
                     \longrightarrow j_*j^*(\ell\boxtimes\ell)).
\]

Only this source must be regular holonomic. On
each diagonal branch its successive terms are
$\ell[1]$ and $\ell$. Tensoring by $C_Q$ replaces
$\ell$ by $L_Q$. The PBW envelope $U_\star L_Q$
and its extraordinary diagonal restriction lie
in the right-module heart. Negative Ext kills
the first term. 

The remaining principal-part
map gives $N\mu_Q$ for an endomorphism $N$.
Inserting the vacuum shows $N=M$. The principal-part
vacuum map is injective on these Hom groups by
the same negative-Ext argument. Interchanging the
slots proves~\eqref{mht:eq:actor-module-detector}.

If $[h_0]|_Y=0$, a relative gauge removes $h_0$
on $Y$. The first actor-weight term then has
one ordinary input and one actor input. Place
them on $Y\setminus\{q\}$ and at $q$, respectively.
The finite Spencer resolution of $\ell$ commutes
with tensoring by arbitrary coefficient complexes.
Algebraic K\"unneth therefore decomposes the
two-output class into a finite sum of bidegrees
$(r,-r)$. 

Since $Q$ lies in degree zero, its
endomorphism complex at the point has only
degree-zero cohomology. After the action, only
$r=0$ survives. The class is thus
$\sum_i M_i\otimes N_i$, with each $M_i$
satisfying~\eqref{mht:eq:actor-module-detector}.

The action comparison identifies this class with
$D\otimes1_Q$. There is no additional zero-column
boundary: maps from the main-diagonal object
$L_Q$ to the open-supported two-root PBW term
vanish by open adjunction. A linear functional
on the finite span of $1_Q,N_1,\ldots,N_r$
expresses $D$ as a linear combination of the
$M_i$. 

Hence $D$ also satisfies
\eqref{mht:eq:actor-module-detector}. But it is
a derivation and kills the vacuum. Subtracting
the derivation identity and inserting that
vacuum gives $D=0$. Faithfulness of localization
on $\mathcal O_X\otimes Q$ contradicts the
chosen nonzero derivation.

The map $f$ has zero pure singleton components.
Restrict to the two-root slice
$(X\setminus\{q\})\times\{q\}$.
The marked form~\eqref{mht:eq:actor-marked-commutator}
gives the binary class of $hf$ as
\[
 [\iota]|_{X\setminus\{q\}}\otimes[h_0]_q
       -[h_0]|_{X\setminus\{q\}}\otimes[\iota]_q,
\]
where $\iota$ is the vacuum pointing of $B$.
The second term is nonzero, since $[h_0]$ is
nonzero there and $[\iota]_q$ acts as $1_Q$.
The two terms have different K\"unneth bidegrees.

This binary class is invariant under arbitrary
pointed homotopies with these endpoints. Write
such a homotopy as $\tau(v)+c(v)\,dv$. Its
singleton terms are $\partial E,\partial F$,
where $E(0)=F(0)=0$ and
$\partial E(1)=\partial F(1)=0$. Off the
collision the source bracket vanishes, and the
binary homotopy equation is
\[
 \tau'_{1,1}=\partial c_{1,1}
              +[\partial E,F']+[\partial F,E'].
\]
The last terms equal
$\partial[E,F']-([E,\partial F])'$.

Integration changes the binary term by a boundary. The endpoint term vanishes. The map through $J$
has zero binary class, so it cannot be homotopic
to $hf$.
\end{proof}

\begin{theorem}\label{mht:thm:ordinary-actor-nonfinality}
Let $Q$ be an ordinary algebraic degree-zero vertex algebra
on $\mathbb A^1_{\mathbb C}$ with a nonzero derivation.
Its full category of factorizing transverse monoid actions
has no final object.
There is a parallel pair of monoid-action maps
that no receiving factorizing monoid action can
equalize up to homotopy, even with unbounded
derived receiving actors.
\end{theorem}

\begin{proof}
In addition to $i:G\to F$, use
$i_0:G\to J\to F$. The action on $G$ is trivial,
so both are morphisms of monoid actions by
Lemma~\ref{mht:lem:relative-derivation-actor}.
Suppose a receiving monoid-action map $r:F\to Z$
equalizes them. Multiplicativity then replaces
the first and third factors of the commutator
word $w$ by units. Its second and fourth factors
cancel in the additive submonoid $H$. Thus
$rw$ is homotopic under $J$ to the map through
$J$.

On the other hand $w=jf$, and $rj$ receives
the nonzero derivation action. Lemma~\ref{mht:lem:geometric-actor-detector}
excludes precisely this homotopy. A final action
would receive $F$ and make the two maps from
$G$ homotopic, contradicting the result.
\end{proof}

For the Heisenberg vertex algebra, the derivation
$D(J)=1$ is nonzero, so the theorem applies.
For $Q=k$, the action category is the slice of
transverse monoids over $J$, and $J$ is final.
These examples concern the stated full actor
category. They do not remove the bounded center
of Theorem~\ref{mht:thm:bounded-geometric-centre},
the ambient endomorphism monoid, or a separately
defined bimodule center.

\section{Self-extensions of the Heisenberg vacuum}

The failure of that universal action leaves a useful
derived endomorphism calculation intact. Let $V_\kappa$
be the Heisenberg vertex algebra at $\kappa\ne0$.
Its modes satisfy
\[
 [J_m,J_n]=\kappa m\delta_{m+n,0},
\]
where the central element acts as the identity.
A module is smooth if every vector is killed by
$J_n$ for all sufficiently large $n$. These are the
modules of the vertex algebra. 

They form a Grothendieck abelian category. Submodules,
quotients and direct sums retain smoothness. Filtered colimits
are exact, and the cyclic smooth quotients form a set of generators.

The vacuum module is
$V_\kappa=k[x_1,x_2,\ldots]$, with
\[
 J_{-n}=x_n,\qquad
 J_n=\kappa n\partial_{x_n}\quad(n\geq1),
 \qquad J_0=0.
\]
For a smooth module $M$, define
\[
 C^\bullet(M)=M\otimes\Lambda(\xi_0,\xi_1,\ldots),
 \qquad dm=\sum_{n\geq0}J_nm\,\xi_n.
\]
Each element uses finitely many exterior generators
and smoothness makes the displayed differential
finite. The nonnegative modes commute, so $d^2=0$.

\begin{proposition}\label{mht:prop:heisenberg-vacuum-ext}
In the smooth module category,
\[
 \operatorname{Ext}^*_{V_\kappa}(V_\kappa,M)
                         =H^*(C^\bullet(M)).
\]
In particular
$\operatorname{Ext}^*(V_\kappa,V_\kappa)=\Lambda(\xi_0)$.
The degree-one class is the nonsplit extension
on which $J_0$ acts by $\varepsilon$ on
$V_\kappa\otimes k[\varepsilon]/(\varepsilon^2)$.
\end{proposition}

\begin{proof}
Let $\mathfrak b$ be the abelian span of $J_n$
for $n\geq0$, and let $\mathfrak a_p$ be the
span for $n\geq p$. Vacuum maps into $M$ are
$M^{\mathfrak b}$. The complex $C^\bullet(M)$
is the union of the finite-mode CE complexes
of $\mathfrak b/\mathfrak a_p$ with coefficients
in $M^{\mathfrak a_p}$.

For a $k[J_0,\ldots,J_{p-1}]$-module $N$, first extend
the action to $\mathfrak b$ by making $\mathfrak a_p$ act
by zero. Induction from this full nonnegative-mode
subalgebra to the Heisenberg algebra at the fixed central
character is exact by PBW. It produces smooth modules:
each vector contains finitely many negative modes, and
sufficiently high positive modes commute past them and
annihilate $N$.
It is left adjoint to $M\mapsto M^{\mathfrak a_p}$.
The latter therefore carries injectives to
injectives. 

For an injective smooth $Y$, each
finite Koszul cochain complex on
$Y^{\mathfrak a_p}$ has no positive cohomology.
Exact filtered colimit gives the same vanishing
for $C^\bullet(Y)$. Since $M\mapsto M\otimes\Lambda$
is exact and its degree-zero cohomology is
$\operatorname{Hom}(V_\kappa,M)$, dimension shifting
proves the Ext comparison.

On the vacuum the differential is
$\sum_{n\geq1}\kappa n\xi_n\partial_{x_n}$.
The degree-minus-one derivation
$\sigma(\xi_n)=x_n/(\kappa n)$ for $n\geq1$,
with $\sigma(\xi_0)=\sigma(x_n)=0$, satisfies
$d\sigma+\sigma d=$ total degree in these
positive-mode pairs. Dividing by this degree
contracts every nonconstant pair monomial.
Only $\Lambda(\xi_0)$ remains.

The stated extension cannot split, because
$J_0$ vanishes on a split sum of vacuum modules
and does not vanish on this extension. Its
class spans the one-dimensional Ext group.
Its square is zero because Ext in degree two
vanishes. This identifies the Yoneda algebra.
\end{proof}

For a Fock module of nonzero charge $\mu$, the
same complex has an additional term
$\mu\xi_0\wedge$. Contraction in $\xi_0$
divided by $\mu$ makes it acyclic. In contrast,
imposing $J_0=0$ as a condition on all modules
removes $\xi_0$ and leaves only the scalars in
the vacuum self-extensions. Thus the answer
depends on the module category.

There is also an explicit coordinate complex
\[
 B_{\mathrm{pt}}=k[u_0,u_1,\ldots]
       \otimes\Lambda(c_0,c_1,\ldots),
 \qquad du_n=\kappa c_{n+1},\quad dc_m=0.
\]
The assignments
$u_n\mapsto n!x_{n+1}$ and
$c_m\mapsto m!\xi_m$ identify it with
$C^\bullet(V_\kappa)$ as a dg algebra. On a
generator both differentials give
$\kappa(n+1)!\xi_{n+1}$.
This point-complex identification retains the
vacuum charge class. Identifying a physical
bulk with this complex on all supports and
collision diagrams requires a further map.
The algebra $\Lambda(\xi_0)$ is not a final
action in the category excluded by
Theorem~\ref{mht:thm:ordinary-actor-nonfinality}.

\section{Complete comparisons and model structures}

A filtered comparison can be tested without
postulating a model category. Let
$f_r:C_r\to D_r$ be a compatible tower of
cochain maps and let
$E_r=\operatorname{Cone}(f_r)$.
Assume the complexes $C$ and $D$ are their
derived inverse limits. Then
\[
 \operatorname{Cone}(f)
       \simeq R\varprojlim_r E_r.
\]
If every $E_r$ is acyclic, so is this derived
limit, and $f$ is a quasi-isomorphism.

Degreewise surjective towers compute their
derived limits by ordinary limits, since
$1-\mathrm{shift}$ is onto on the product
complex. These are the explicit hypotheses
used in the continuous relative comparison.

A stronger constructive criterion gives a
homotopy. If $h_r:E_r\to E_r[-1]$ commute
with the transitions and satisfy
$dh_r+h_rd=1$, then $\varprojlim h_r$
contracts the ordinary limit. Contractions
chosen independently at each finite stage
do not define that limit homotopy.

Completeness matters even for zero
differentials. For nonzero finite-dimensional
spaces $V_r$, the inclusion
\[
 \bigoplus_{r\geq1}V_r\longrightarrow\prod_{r\geq1}V_r
\]
is an isomorphism on every finite tail
quotient, but its ordinary cokernel is
nonzero. The direct sum is not complete for
the tail filtration. Thus this example
does not contradict the derived-complete
comparison criterion.

A Quillen presentation adds hypotheses to
this calculation. In ordinary nonnegative
homological complexes, let $S^n$ be $k$ in degree $n$ and
let $D^n$ have $k$ in degrees $n,n-1$ with
identity differential, for $n\geq1$.
The generating sets are
\[
 I=\{0\to S^0\}\cup
       \{S^{n-1}\to D^n:n\geq1\},\qquad
 J=\{0\to D^n:n\geq1\}.
\]
The degree-zero sphere belongs to $I$,
and $0\to S^0$ is not acyclic. A filtered
version must specify its filtration-cell
functors as well as smallness, recognition
and tensor compatibility. Exact inverse
limits alone do not establish these
model-category axioms.

Once actual algebra and coalgebra model
categories and a Quillen bar--cobar
adjunction have been constructed, a
restricted equivalence has a precise
criterion. Choose full derived
subcategories preserved by the two
derived functors. If their derived unit
and counit are equivalences on those
subcategories, the functors are inverse
there: their unit and counit are the
required natural equivalences. An
equivalence of the whole categories
requires these hypotheses on every
relevant object. A length filtration or
a finite quotient does not supply the
missing global unit or counit.

\section{The open--closed comparison}

Let $\mathcal B$ be the independently
defined bulk observables and let
$\mathcal C_\partial$ be the chosen
category of boundary conditions. Suppose
its complete factorization-bimodule
category and diagonal have been
constructed. A coherent bulk action is
then a map to
\[
 R\operatorname{Hom}_{\mathrm{Bimod}^{\mathrm{fact}}
                          (\mathcal C_\partial)}
            (\Delta_\partial,\Delta_\partial),
\]
with its support, coefficient base,
completion and operation structure
specified. Its evaluation at one brane
is a further map to that brane's
Hochschild cochains.

At the Lie level the first equations
are explicit. Let $\mathfrak g$ be a dg Lie algebra of closed coefficients
and $A$ a boundary dg algebra. Put
$\mathfrak h=C^\bullet(A,A)[1]$, with its compatible
shifted differential and Gerstenhaber bracket. A degree-one
twisting cochain
\[
 \tau:\overline C_*^{\mathrm{CE}}(\mathfrak g)
                                      \longrightarrow\mathfrak h
\]
satisfies
\begin{equation}\label{mht:eq:centre-action-mc}
 d\tau+\frac12[\tau,\tau]=0.
\end{equation}

The bracket uses the CE coproduct and
the bracket of $\mathfrak h$. The first
Taylor coefficient is a chain map. The second supplies its bracket
homotopy. Later coefficients fill the
higher action equations. 

This is an
$L_\infty$ action in Hochschild
cochains. A mixed geometric action also
needs compatibility with the braces,
units, disjoint insertions and all
collision maps.

These additional identities can be
treated by an operadic convolution
complex once a resolution of the actual
colored operation operad is supplied.
If that resolution is a cobar
$\Omega\mathcal D$, its convolution
with the endomorphism operations of
$(\mathcal B,A)$ has Maurer--Cartan
elements precisely the corresponding
homotopy-coherent structures. Twisting
by the fixed bulk and boundary
operations leaves the relative action
problem. Merely listing desired
operation-preserving maps does not
define this deformation complex.

For any such complete filtered
convolution complex, a partial solution
has a closed curvature in the first
unsolved filtration degree. A primitive
kills that curvature and extends the
solution. Completeness and compatible
choices produce a full action. The
complex and its filtration must be
constructed on the chosen supports
before their obstruction classes are
defined.

Even a completed action need not be an
equivalence. After the map exists,
its comparison cone is the remaining
test. On derived complete coefficient
towers, the preceding cone criterion
provides a sufficient proof. At the
category level, evaluation at a brane
must also be related to the diagonal
by the actual bimodule functor.

\begin{problem}\label{mht:qn:physical-centre-comparison}
Construct a bulk action for the mixed Hamiltonian theory
on the declared boundary category. Retain its coefficient
topology, both units, supported collision operations and
their higher homotopies. Determine
whether its comparison cone is
acyclic. If a restricted actor or
brane category is required, construct
its relation to the physical boundary
conditions and identify which
operations it preserves.
\end{problem}

The finite coordinate identities and
continuous relative center give
explicit algebraic domains for this
problem. The geometric obstruction
rules out a particular unrestricted
final-action formulation on a curve.
Neither determines the missing mixed
bulk action. Quantization and
reconstruction must therefore retain
the chosen recipient and its
comparison maps.

\section{A normalized junction on a curve}
\label{mht:sec:ramond-junction}

An action of two boundary algebras on one complex does not by
itself reconstruct either boundary category. The obstruction can
already occur in a free theory on a holomorphic curve. We shall
exhibit nonzero modules invisible to the junction, then an
extension on which its four generating fields cannot act
compatibly. Finally, a localization will distinguish polynomial
local operators from derived mapping objects.

\subsection{States, fields and mapping complexes}

Put $R=\mathbb C[[q]]$ and $S_\ell=R/(q^\ell)$, where $\ell\geq1$.
The parameter $q$ has degree, weight and energy zero.
On a marked spin coordinate define the Ramond complex
\begin{align*}
 P_\ell&=S_\ell[x_i,y_n]\otimes\Lambda(\theta_i,\eta_n),
 \qquad i\geq0,\quad n\geq1,\\
 D&=\sum_{i\geq0}x_i\partial_{\theta_i}
       -\sum_{n\geq1}\eta_n\partial_{y_n}.
\end{align*}

The cohomological degrees of $x,y,\theta,\eta$ are respectively
$0,0,-1,1$. Give $x_i,\theta_i$ weight and energy
$(1,2i+1)$, and $y_n,\eta_n$ weight and energy $(-1,2n-1)$.
Every state is a finite polynomial. The complete complex
$P=\varprojlim_\ell P_\ell$ allows a different finite polynomial
at each coefficient of $q$. It does not take products over
energy at a fixed coefficient.

The even boundary algebra $A_0$ is the commutative differential
polynomial algebra generated by $\gamma$. The odd boundary
algebra $A_1$ is its exterior counterpart generated by
$\widehat c$. Both have zero differential and zero
self-contractions. Their generating fields on $P_\ell$ are
\begin{align*}
 \Gamma(z)&=\sum_{i\geq0}x_i z^{i-1/2}
                   -\sum_{n\geq1}\partial_{y_n}z^{-n-1/2},\\
 C(z)&=\sum_{n\geq1}\eta_n z^{n-1/2}
                   +\sum_{i\geq0}\partial_{\theta_i}z^{-i-1/2}.
\end{align*}

The two additional fields
\begin{align*}
 \widehat\beta(z)&=\sum_{n\geq1}y_n z^{n-1/2}
                   +\sum_{i\geq0}\partial_{x_i}z^{-i-1/2},\\
 B(z)&=\sum_{i\geq0}\theta_i z^{i-1/2}
                   +\sum_{n\geq1}\partial_{\eta_n}z^{-n-1/2}
\end{align*}
satisfy, with graded commutators,
\[
 [D,B]=\Gamma,\qquad [D,\widehat\beta]=-C,
 \qquad [D,\Gamma]=[D,C]=0.
\]
Writing
$\Delta(u,z)=\sum_{r\in\mathbb Z}u^{r-1/2}z^{-r-1/2}$,
the only nonzero brackets between the four generating fields
are
\[
 [\widehat\beta(u),\Gamma(z)]=\Delta(u,z),\qquad
 [C(u),B(z)]=\Delta(u,z).
\]

These identities follow coefficientwise from the Weyl and
exterior contraction relations. We call the four fields with
these relations the \emph{normalized quartet}. The physical
odd field $c=qC$ is a separate normalization. No construction
here inverts $q$.

Translation is the operator
\begin{align*}
 T={}&\sum_{i\geq0}(i+\tfrac12)
       (x_{i+1}\partial_{x_i}+\theta_{i+1}\partial_{\theta_i})\\
   &+\sum_{n\geq1}(n+\tfrac12)
       (y_{n+1}\partial_{y_n}+\eta_{n+1}\partial_{\eta_n})
       -\tfrac12y_1x_0+\tfrac12\eta_1\theta_0.
\end{align*}
It has weight zero and energy two, commutes with $D$, and obeys
$[T,F(z)]=\partial_zF(z)$ for each quartet field. The two
multiplication terms supply the change from creation to
annihilation modes. In particular,
\[
 T1=-\tfrac12y_1x_0+\tfrac12\eta_1\theta_0
        =D(-\tfrac12y_1\theta_0).
\]

The constant state is closed but is not translation invariant.
The boundary vacuum field acts as the identity. This unit
condition does not turn the constant junction state into the
boundary vacuum.

To specify the recipients, write
$X_i=x_i$, $G_n=-\partial_{y_n}$ for the even boundary modes
and $H_n=\eta_n$, $K_i=\partial_{\theta_i}$ for the odd ones.
Let $\mathcal R_{0,\ell}$ be the category of bigraded dg modules
for the commuting $X_i,G_n$. Require the following support
condition: for every vector $v$, some $N,h$ satisfy
\[
 J_{N,h}v=0,\qquad
 J_{N,h}=(G_n:n>N)+(G_1,\ldots,G_N)^h.
\]

Let $\mathcal R_{1,\ell}$ be the corresponding category for the
anticommuting $H_n,K_i$, where every vector is killed by all
sufficiently high $K_i$. All modules are unital.

Differential
polynomials act by products of derivatives of the generating
field. Move annihilators to the right: their indices and powers
are bounded on an input, and a fixed Laurent coefficient then
bounds the remaining creation indices. Thus these products
are fields with finite coefficients on each input.

Morphisms in these categories strictly intertwine the boundary
modes, with the cohomological signs. The restricted Hom complex
is the direct sum of the Hom complexes of fixed weight and
energy shift. Its derived version is formed in the stated
category at each $S_\ell$, before taking the derived inverse
limit in $\ell$. One may also retain a translation operator
with the displayed mode commutators and require morphisms to
commute with it. These mapping complexes do not impose
locality with the quartet. Local
operator fields will be considered separately below.

\subsection{Nonzero modules invisible to the junction}

Define two single-boundary modules, each with zero differential,
by
\[
 N_{0,\ell}=S_\ell[x_i,y_n],\qquad
 N_{1,\ell}=\Lambda_{S_\ell}(\theta_i,\eta_n).
\]
Use the boundary actions just given. Their translations $T_0$
and $T_1$ are respectively the even and odd parts of $T$,
including the terms $-y_1x_0/2$ and $\eta_1\theta_0/2$.
Both modules obey the support conditions: polynomial degree
bounds the even annihilators, and a finite exterior monomial
has only finitely many $\theta$ indices. Their energies are
nonnegative and each energy piece is finite free over
$S_\ell$. The constant vector gives a nonzero cohomology class.

\begin{theorem}\label{mht:thm:ramond-orthogonals}
For $a=0,1$, the restricted derived mapping complex from
$P_\ell$ to $N_{a,\ell}$ vanishes in $\mathcal R_{a,\ell}$,
both with and without translation. The vanishing persists
after the prescribed derived coefficient completion.
Consequently the junction does not generate either category,
even in the full derived subcategory whose objects have
nonnegative energy and finite free energy pieces.
\end{theorem}

\begin{proof}
Put $B_0=S_\ell[x_i]$ and $B_1=\Lambda_{S_\ell}(\eta_n)$.
The even support condition identifies its recipient with
comodules in bigraded $B_0$-modules for the coalgebra
\[
 C_+=\bigoplus_\alpha S_\ell e_\alpha,
 \qquad \Delta e_\alpha=\sum_{\beta+\gamma=\alpha}
                         e_\beta\otimes e_\gamma.
\]
Here $\alpha$ ranges over finite multiindices. The coaction
of $v$ is $\sum_\alpha e_\alpha\otimes G^\alpha v$.
The support condition makes it finite.

The odd recipient is
similarly the category of comodules for the exterior
coalgebra $\Lambda(\theta_i)$ in bigraded $B_1$-modules.
The field construction above identifies these mode
presentations with the respective boundary actions.

In either presentation the forgetful functor $U$ is exact.
Its right adjoint $F(V)=C\otimes_{S_\ell}V$ is exact because
the coalgebra is a direct sum of free coefficient modules.
On the even side identify $e_\alpha$ with
$(-y)^\alpha/\alpha!$. On the odd side use the exterior
cofree module. Thus $N_a=F(B_a)$.

If $I$ is a K-injective
base-module resolution, $FI$ is K-injective: Hom from an
acyclic comodule is Hom from its acyclic underlying module
to $I$. This proves the unbounded derived adjunction
\[
 R\operatorname{Hom}_{\mathcal R_a}(M,FV)
       \simeq R\operatorname{Hom}_{B_a}(UM,V)
\]
in every internal shift. Direct sums of those shifts are
exact.

The polynomial Euler contractions of the pairs
$(\theta,x)$ and $(y,\eta)$ give $UP_\ell\simeq S_\ell$
as a $B_a$-module. It remains to show
$R\operatorname{Hom}_{B_a}(S_\ell,B_a)=0$ in each shift.

For $B_0$, use the semifree Koszul resolution
$B_0\otimes\Lambda(\theta_i)$ with $d\theta_i=x_i$.
A degree-$p$ Hom cochain is an alternating family
$(f_I)_{|I|=p}$ of finite polynomials. Its index support may
be infinite. If $\epsilon_i$ denotes a formal dual exterior
symbol, the raw differential is
$(-1)^{p+1}\sum_i x_i\epsilon_i\wedge-$.
Multiplication of degree $p$ by $(-1)^{p(p+1)/2}$ changes
this to $\delta=\sum_i x_i\epsilon_i\wedge-$.
Every output coefficient of this formula is a finite sum.

Define
\[
 Q_jf=\frac{f-f|_{x_j=0}}{x_j},\qquad
 h_j=\iota_{\epsilon_j}Q_j.
\]
Let $\pi_j$ retain the terms containing $\epsilon_j$ and
set $x_j=0$ in their coefficients. Division in $Q_j$ is
polynomial division of an expression divisible by $x_j$.
On terms with and without $\epsilon_j$, the two contraction
identities give
\[
 \delta h_j+h_j\delta=1-\pi_j.
\]

The other indices anticommute with $h_j$, and each $\pi_j$
commutes with $\delta$. For any $m$ distinct consecutive
indices form
\[
 H_m=h_0+\pi_0h_1+\cdots+\pi_0\cdots\pi_{m-2}h_{m-1}.
\]
Then
\[
 \delta H_m+H_m\delta=1-\pi_0\cdots\pi_{m-1}.
\]

The last projection requires all $m$ exterior symbols, so
it kills a cochain of degree $p<m$. A degree-$p$ cocycle
therefore has the primitive $H_{p+1}f$. Negative degrees
are zero. Each primitive uses finitely many queries of the
family and produces a finite polynomial on every input.
Removing $\epsilon_j$ cancels the internal degree of division
by $x_j$, so the primitive preserves the prescribed shift.

For $B_1$, take the semifree resolution
\[
 B_1\otimes S_\ell[s_1,s_2,\ldots],
 \qquad ds_n=\eta_n,\qquad |s_n|=0,
\]
using divided monomials and the same internal degrees for
$s_n$ and $\eta_n$. A degree-$p$ Hom cochain is a family
$F_\alpha\in B_1^p$. Its raw differential is
$-\sum_{n:\alpha_n>0}\eta_nF_{\alpha-e_n}$: the sign from
graded left-linearity cancels the degree-dependent Hom
sign. Multiplication of degree $p$ by $(-1)^p$ gives
\[
 (\delta F)_\alpha=\sum_{n:\alpha_n>0}
                         \eta_nF_{\alpha-e_n}.
\]

Set
\[
 (h_jF)_\alpha=\iota_{\eta_j}F_{\alpha+e_j},\qquad
 (\pi_jF)_\alpha=
 \begin{cases}
 \eta_j\wedge\iota_{\eta_j}F_\alpha,&\alpha_j=0,\\
 0,&\alpha_j>0.
 \end{cases}
\]

Exterior contraction again gives
$\delta h_j+h_j\delta=1-\pi_j$.
The same telescoping sum, now over $j=1,\ldots,m$, has a
final projection whose output contains $m$ distinct
$\eta$'s. It vanishes for $m=p+1$. The resulting primitive
has finite outputs and unchanged internal shift.

The
Euler contractions prove that the two resolutions resolve
$S_\ell$. Their filtrations by the number of free generators
prove semifreeness, hence K-projectivity. No bound on the
cohomological degree of the target was used.

It remains to retain translation. Let $\delta_T$ be the
derivation of the mode algebra given by commutation with
$T$. Induction to its Ore extension has underlying module
$S_\ell[T]\otimes M$, with
\[
 a(T^r\otimes v)=\sum_{k=0}^r(-1)^k\binom rk
                 T^{r-k}\otimes\delta_T^k(a)v.
\]
It is exact and preserves the support condition.
Indeed, $J_{N,h}v=0$ implies
$J_{N+r,h+r}(T^r\otimes v)=0$: at most $r$ positive modes
are differentiated in any summand, leaving at least $h$
undifferentiated factors. High odd annihilators similarly
increase their cutoff by at most $r$.

Every translated module has the two-term relative induced
resolution whose map sends
$T^r\otimes v$ to $T^{r+1}\otimes v-T^r\otimes T_Mv$.
Its injectivity follows from the leading power of $T$.
Its cokernel is $M$. The two induced terms give acyclic
derived Hom complexes by the translation-forgetting result.
Their cone is therefore acyclic. This argument does not
require the earlier primitives to commute with translation.

Finally the derived inverse limit of the zero mapping
objects is zero. The targets remain nonzero, which proves
failure of generation.
\end{proof}

For finitely many Koszul pairs, the last exterior class
survives in top degree. Infinite vanishing uses a longer
finite set of indices in each fixed degree. It is not a
limit of assertions that the finite complexes are acyclic.
If a tensor--Hom adjunction for the boundary junction has
been constructed, its counit on $N_a$ has source zero and
nonzero target. This prevents an equivalence on the stated
recipient. It says nothing by itself about dualizability
in a different category with a different fusion tensor.

A smaller recipient can have a positive generation theorem.
For $A=S_\ell[x]$, the complex
$[A\theta\xrightarrow{x}A]$ in degrees $-1,0$ is a compact
nongenerator of $D(A)$: Hom into $A[x^{-1}]$ is contracted
by $-x^{-1}$. It generates the independently defined
subcategory of complexes whose localization at $x$ vanishes.
Indeed the local-cohomology complex
$[A\to A[x^{-1}]]$ is the filtered colimit of
$(A/x^n)[-1]$. Each $A/x^n$ is built by extensions from
$A/x$, and tensoring this local-cohomology complex with
an $x$-torsion object recovers that object. A restriction
of the recipient can therefore repair generation, but its
support condition must be specified independently.

\subsection{A stable extension without a compatible quartet}

Requiring every object to admit the four fields is another
possible restriction. If morphisms still intertwine only
one boundary, that restriction fails to be stable. The
failure is detected by two commuting even boundary modes.

Define the operator
\[
 \Omega=\sum_{i\geq0}(2i+1)\eta_{i+1}\partial_{\theta_i}
            =-\operatorname{Res}_z C(z)\partial_zC(z).
\]
It has cohomological degree two, weight $-2$ and energy
zero. The sum is finite on each state. It commutes with
$D$, all even boundary modes, and $T$. For translation,
the coefficient of each adjacent pair cancels by
\[
 (2i+1)(i+\tfrac32)-(2i+3)(i+\tfrac12)=0;
\]
the remaining multiplication term contains $\eta_1^2=0$.

Adjoin an odd variable $\varepsilon$ of cohomological degree
$-1$, weight two and energy zero. For $a\in S_\ell$ put
\[
 E_a=P_\ell\otimes\Lambda(\varepsilon),\qquad
 d_a=D+a\varepsilon\Omega,
\]
where $D\varepsilon=T\varepsilon=0$ and the even boundary
actions are unchanged. Closedness of $\Omega$ gives
$d_a^2=0$ and $[T,d_a]=0$. The exact sequence
$\varepsilon P_\ell\to E_a\to P_\ell$ places $E_a$ in
the stable subcategory generated by $P_\ell$ and its
internal shifts. All its energy pieces are finite free.

The change of polynomial coordinates
$x'_i=x_i+a(2i+1)\varepsilon\eta_{i+1}$ gives underlying
cohomology $\Lambda_{S_\ell}(\varepsilon)$. It does not
preserve the boundary action.

\begin{proposition}\label{mht:prop:ramond-stable-cone}
If $a\ne0$, the even boundary module $E_a$ is not
quasi-isomorphic to one admitting a compatible normalized
quartet. Let $\mathcal S$ be a full stable subcategory
of the even boundary derived category, closed under
equivalence and containing $P_\ell$ and its internal shifts.
Some object of $\mathcal S$ then admits no compatible
quartet, even up to quasi-isomorphism.
\end{proposition}

\begin{proof}
Retain $x=X_0$ and $g=G_1=-\partial_{y_1}$, and take the
derived specialization $g=0$. The two-term free resolution
of $S_\ell$ over $S_\ell[g]$ computes this specialization.
Since $g$ is degreewise surjective, its cone is, up to the
common cohomological shift, its kernel. This is a statement
about complexes. Integration with respect to $y_1$ need
not be a chain map. The kernel omits $y_1$.

Contract the remaining pairs after the coordinate change
above. The residual complex is
\begin{gather*}
 K_a=S_\ell[x]\otimes\Lambda(\theta,\varepsilon,\eta),\\
 d\theta=x+a\varepsilon\eta,\qquad d\varepsilon=d\eta=0.
\end{gather*}
Division by the monic polynomial $x+a\varepsilon\eta$
gives its cohomology and the action of $x$:
\[
 H(K_a)=\Lambda_{S_\ell}(\varepsilon,\eta),\qquad
 x[1]=-a[\varepsilon\eta].
\]
The cohomology is free over $S_\ell$, so this action is
nonzero for every nonzero $a$, including nilpotent $a$.

For example, if $u=a\varepsilon\eta$, then $u^2=0$ and
\[
 \frac{x^m-(-u)^m}{x+u}=x^{m-1}-ux^{m-2}\quad(m\geq2).
\]
Multiplying this quotient by $\theta$ gives the explicit
contraction onto evaluation at $x=-u$. The cases $m=0,1$
are immediate. No polynomial-degree quotient is involved.

In a compatible quartet the coefficient $b_0$ of $B$
commutes with every even boundary mode and satisfies
$[d,b_0]=x$. It therefore gives a null-homotopy of $x$
after derived specialization at $g=0$. The nonzero action
just calculated contradicts this. Restriction and derived
specialization preserve quasi-isomorphisms, so the
contradiction survives every replacement in the even
boundary derived category. Stability would include $E_a$,
which proves the final assertion.
\end{proof}

The same extension occurs without dividing the physical
odd field by $q$. The expression
$-\operatorname{Res}c\partial c=q^2\Omega$ gives $a=q^2$.
It is nonzero over $S_\ell$ for $\ell\geq3$ and over $R$.
The lower two precisions do not detect it. Thus an
additional support condition satisfied by $P$ cannot
remove this extension while preserving full stable
closure and the original morphisms.

One can restore the differential identity alone by changing
the odd field to
\[
 B_a=B-2a\varepsilon\partial_z\widehat\beta.
\]
It obeys $[d_a,B_a]=\Gamma$ and translation covariance.
Its mixed bracket with $\Gamma$ is
$2a\varepsilon\partial_z\Delta$, however, so it is not
the original quartet. At the two selected modes,
$[g,\theta_0-a\varepsilon y_1]=a\varepsilon$.
The defect occurs in the boundary-linearity required
of the homotopy, not in the underlying scalar contraction.

There is also a positive category with stricter morphisms.
Let $\mathcal W_\ell$ be the dg Weyl--Clifford algebra of
finite polynomial operators in all quartet modes. In
normal order its symbols are the polynomial generators
$x,y,\xi,\zeta$ and exterior generators
$\theta,\eta,\kappa,\pi$, where
$\xi=\partial_x$, $\zeta=\partial_y$,
$\kappa=\partial_\theta$ and $\pi=\partial_\eta$.
The differential pairs them by
\[
 d\theta=x,\qquad dy=-\eta,
 \qquad d\xi=-\kappa,
 \qquad d\pi=-\zeta.
\]

On normal symbols, the reverse pairing divided by positive
total polynomial degree is an Euler contraction. Every
symbol uses finitely many modes, so the contraction proves
$S_\ell\to\mathcal W_\ell$ is a quasi-isomorphism of dg
algebras. Evaluation $w\mapsto w1$ is a quasi-isomorphism
$\mathcal W_\ell\to P_\ell$ of left dg modules. Hence
\[
 D(\mathcal W_\ell)\simeq D(S_\ell),\qquad
 R\operatorname{End}_{\mathcal W_\ell}(P_\ell)\simeq S_\ell,
\]
and $P_\ell$ is a compact generator in this ordinary
module category.

Its endomorphisms recover scalars, not
the opposite boundary algebra. The category has changed
its morphisms, and no boundary fusion has been constructed.
The primitive diagonal action on a tensor product of two
Weyl modules has $[\partial,x]=2$. It does not define a
tensor for the relation $[\partial,x]=1$.

\subsection{Polynomial local operators and derived descent}

The preceding derived Hom complexes consist of boundary
intertwiners. A different prescription uses fields local
with the quartet. At each coefficient of $q$, use normally
ordered polynomial expressions in the original creation
and annihilation symbols, with bounded total degree.
Mode indices may range without bound. Require every
Laurent coefficient to act finitely on each input and
the field to be lower truncated on that input.

Impose strict linearity for one boundary and locality with
each quartet field. Translation covariance may also
be imposed. Denote this complex of local operator fields
from $M$ to $N$ by $\mathcal A_{\mathrm{pol}}(M,N)$,
with differential the commutator with $D$.

Locality means that a power of $u-z$ kills each graded
commutator. It does not mean that every coefficient is
an intertwiner of all quartet modes. The polynomial
degree restriction is part of the definition of
$\mathcal A_{\mathrm{pol}}$. Its derived invariance
must therefore be checked in that precise class.
A derived bifunctor is called reduced here if it
vanishes whenever either argument is zero.

\begin{proposition}\label{mht:prop:ramond-local-descent}
Put $Q_\ell=P_\ell[x_0^{-1}]$, with its extended quartet
and translation. Its underlying complex is acyclic,
but $\mathcal A_{\mathrm{pol}}(Q_\ell,Q_\ell)$ has a
nonzero identity class. The rectangular complex
$\mathcal A_{\mathrm{pol}}(P_\ell,Q_\ell)$ likewise has
a nonzero inclusion class. Thus this prescription does
not descend, objectwise or on the displayed pair, to a
reduced bifunctor on the localization at underlying
quasi-isomorphisms.
\end{proposition}

\begin{proof}
Localization leaves all polynomial-input support bounds
intact. The fields and translation extend by the usual
derivative of $x_0^{-1}$. Multiplication by
$h=\theta_0/x_0$ satisfies
$Dh+hD=1$, so $Q_\ell\to0$ is an underlying
quasi-isomorphism. This contraction need not intertwine
the chosen boundary or commute with translation.

Every operator in the polynomial class preserves
$P_\ell\subset Q_\ell$. Conversely, such an operator
on $P_\ell$ extends uniquely to $Q_\ell$. Its
$\partial_{x_0}$-order is bounded, and a finite input
involves only finitely many other derivative indices.
Normal coefficients can be recovered by evaluating on
monomials of increasing degree. The diagonal entries
are factorials, which are units in $S_\ell$.
This also proves faithfulness after localization.

Apply that argument to each Laurent coefficient. All
identities for the differential, boundary-linearity,
locality and translation are identities of the same
finite-degree normal operators. They hold on $Q_\ell$
exactly when they hold on $P_\ell$. The polynomial
diagonal and rectangular complexes are therefore
identified with the corresponding complex on
$P_\ell$, with identity or inclusion corresponding to
the identity operator.

Let $\epsilon_P:P_\ell\to S_\ell$ be the constant-term
functional. It satisfies $\epsilon_PD=0$ and $D1=0$.
If the identity were $[D,F]$ for a degree-$-1$
polynomial field, evaluation on $1$ would give
$D(F1)=1$. Applying $\epsilon_P$ contradicts this
coefficientwise.

Thus the identity and inclusion
classes are nonzero. A reduced derived bifunctor
must vanish when one argument is the zero object.
Thus $Q_\ell\simeq0$ proves the claimed failure.
\end{proof}

The same argument holds after coefficientwise completion.
A polynomial primitive over $R$ would reduce to an
impossible primitive modulo $q$.

This example allows negative energy, with an infinite
energy-zero piece. It does not rule out a construction
on a smaller positive-energy category, nor a derived
replacement agreeing with the polynomial prescription
only at $P$. Both require their own mapping objects
and comparison maps.

On the even boundary, enlarging the coefficient class
does produce a local contraction. Write
\[
 a(z)=\sum_{i\geq0}x_i z^i,\qquad
 p(z)=\sum_{n\geq1}\partial_{y_n}z^{-n},\qquad
 \Gamma=z^{-1/2}(a-p).
\]
Since $x_0$ is invertible on $Q_\ell$, the formal series
$a(z)^{-1}$ is defined. The operator field
\[
 I(z)=z^{1/2}\sum_{r\geq0}a(z)^{-r-1}p(z)^r
\]
is an inverse of $\Gamma(z)$. For a fixed input the
$r$-sum ends at its total $y$-degree. Each remaining
Laurent coefficient involves finitely many variables.

Consequently
\[
 H(z)=B(z)I(z),\qquad [D,H(z)]=1
\]
is a well-defined lower-truncated field. It commutes
with $\Gamma$, is translation covariant, and is local
with the quartet. Indeed
\begin{align*}
 [\widehat\beta(u),H(z)]&=-\Delta(u,z)B(z)I(z)^2,\\
 [C(u),H(z)]&=\Delta(u,z)I(z),
\end{align*}
while its brackets with $\Gamma$ and $B$ vanish.
Multiplication by $u-z$ kills the displayed distributions.

This inverse field requires unbounded orders of poles
in $x_0$ across its Laurent coefficients. After setting
$x_i=0$ for $i\geq2$, its $\theta_0$ component on the
constant state is
\[
 \frac{\theta_0}{x_0+x_1z}
   =\sum_{m\geq0}(-1)^m\theta_0x_1^m x_0^{-m-1}z^m.
\]
It therefore lies outside the polynomial class.

Allowing a fixed finite pole order does not supply
the same construction. In fact, locality with
$\widehat\beta$ rules out a primitive in the class
of finite-degree normal operators with a globally
bounded Laurent pole order in $x_0$. Let $K$ be a
locality order at a fixed coefficient of $q$. Multiplication
by $(u-z)^K$ gives the $K$-th finite-difference equation
in the annihilator index. For a normal symbol $\Sigma$
its solutions have the form
\[
 \partial_{x_i}\Sigma
 =\sum_{k=0}^{K-1}G_k(z)(i-\tfrac12)^{\underline{k}}
                              z^{i-1/2-k}.
\]

Here $t^{\underline{k}}=t(t-1)\cdots(t-k+1)$, and
$t^{\underline{0}}=1$. The falling factorials form a
basis of polynomial sequences of degree less than $K$.
Evaluation at $i=1,\ldots,K$ has an invertible Vandermonde
matrix. It therefore expresses the row $i=0$ as a
Laurent scalar combination of those rows. Hence
\[
 \left(\partial_{x_0}-\sum_{i=1}^K
                    c_i(z)\partial_{x_i}\right)\Sigma=0
\]
for suitable $c_i(z)$. For $K=0$ use $\partial_{x_0}$.

If $\Sigma$ has largest pole $x_0^{-m}$ with coefficient
$\Sigma_{-m}$ independent of $x_0$, the last equation
has coefficient $-m\Sigma_{-m}$ at $x_0^{-m-1}$.
The other derivatives cannot increase that pole order.
Positive integers are units in $S_\ell$, so this is
impossible. No pole can occur, and the polynomial
augmentation argument applies.

Strict odd boundary-linearity has a stronger obstruction.
For a degree-$-1$ field $F$ commuting with all $K_i$,
the vector $F(z)1$ lies in every $\ker\partial_{\theta_i}$.
It is therefore $\theta$-free. Such a vector has
nonnegative cohomological degree, even with Laurent
coefficients in the even variables, and so $F(z)1=0$.
It cannot satisfy $[D,F]=1$. The nonzero bracket
$[C,H]=\Delta I$ is exactly where the positive
inverse-field construction leaves this strict odd
boundary class.

The normalized curve model thus separates three
questions. Strict derived intertwiners need not detect
all boundary objects. Objects carrying a quartet need
not remain so after cones of boundary-linear maps.
Polynomial local fields need not define derived
mapping objects.

Requiring all quartet modes to be
intertwined repairs generation in an ordinary scalar
module category. Allowing inverse fields repairs the
even local contraction on a larger coefficient class.
Neither construction identifies the physical boundary
fusion, its unit, or its supported collision operations.
Those comparisons remain part of
Problem~\ref{mht:qn:physical-centre-comparison}.
